% ---------------------------------------------------------------
%  AUTO-GENERATED by tools/build_compendium.py — DO NOT EDIT.
%  Edit the individual reading notes and re-run the builder.
% ---------------------------------------------------------------
\documentclass[11pt]{book}
\usepackage[margin=1in]{geometry}
\usepackage{amsmath}
\usepackage{enumitem}
\usepackage{graphicx}
\usepackage{parskip}
\usepackage{titlesec}
\usepackage{xcolor}
\usepackage[hidelinks]{hyperref}

% Sources are written as articles, so their deepest heading is \subsection.
% Under \part that lands at \subsubsection — keep it numbered and in the ToC.
\setcounter{secnumdepth}{3}
\setcounter{tocdepth}{2}

\title{\textbf{2026 Comprehensive Exam}\\[0.3em]
\large Collected Guides and Notes}
\author{Ryan Gibbons}
\date{\today}

\begin{document}
\maketitle
\tableofcontents


% ============================================================
\part{Anderson 1983}
\begin{center}
Anderson, Benedict. \textit{Imagined Communities: Reflections on the Origin and Spread of Nationalism}. New York: Verso, 1983.
\end{center}

\section{My Notes}

\subsection{Main Idea}



%--------------------------------------------------
\subsection{General Notes}
%--------------------------------------------------

\begin{itemize}[leftmargin=*, itemsep=4pt]
    \item Nations are inherently ``imagined'' or otherwise products of conception rather than explicit lines: most members of a nation will never meet most of the other members, but still a sense of commonality exists. The nation must be limited, sovereign, and considered as a community; these align with four portions of the definition of the nation:
    \begin{enumerate}
        \item Imagined (not the same as fake -- just that it cannot be readily conceived in its entirety as a face-to-face community would)
        \item Limited / Exclusive
        \item sovereign
        \item Exists as a community
    \end{enumerate}
    \item
\end{itemize}


\subsection{Important Specific Factual Claims}

\begin{itemize}[leftmargin=*, itemsep=4pt]
    \item Nationalism exists on a separate dimension from fascism, liberalism, etc. -- emerged from the decline of a religious / dynastic justification for peoplehood that declined in the Enlightenment period and continued therafter
    \item A great deal of nation-building was facilitated by the development of effective communication technology (i.e. printing press) and capitalistic markets, which facilitated exchange particularly by means of standardized (and particularized) language
    \item Colonial empires developed strategic nationalistic characteristics primarily to protect dynastic and aristocracies by creating a common vision of peoplehood
    \item Origins largely from Latin American Creole cultures
    \item Three primary institutions which supported this construction particularly within colonial societies:
    \begin{enumerate}
        \item Censuses: allowed for distinct, mutually exclusive counts and categorizations
        \item Maps: allowed for distinct borders and the creation of commonly understood delineations
        \item Museums: display of stories of peoplehood and the origins of the nation
    \end{enumerate}
\end{itemize}


\subsection{Connections}

\begin{itemize}[leftmargin=*, itemsep=4pt]
    \item Thinking back to POLS 111 and stories of peoplehood -- challenges, strengths, triumps, common enemies, common goals
    \item Cites Weber's discussions of legitimate authority and the impacts of charismatic leadership
\end{itemize}


\subsection{My Critiques}

\begin{itemize}[leftmargin=*, itemsep=4pt]
    \item
\end{itemize}


% ============================================================
\part{Boix 2003}
\begin{center}
Boix, Carles. \textit{Democracy and Redistribution}. New York: Cambridge University Press, 2003.
\end{center}

\section{My Notes}

\subsection{Main Idea}



%--------------------------------------------------
\subsection{General Notes}
%--------------------------------------------------

\begin{itemize}[leftmargin=*, itemsep=4pt]
    \item Democracy is a redistributive activity and therefore stands to prevail when the impacts of this redistribution can be mitigated. Two pathways through which this can occur:
    \begin{enumerate}
        \item Economic equality: there is little change to be expected from ``fair'' redistribution, so establishing a redistributive institution is unlikely to be particularly disruptive
        \item Capital mobility: the friction to redistribution is low/lessened such that the redistribution can occur without broad societal upheaval
    \end{enumerate}
    \item If capital is significantly mobile such that it could be redirected abroad in the face of threats/abuse, then the democratic state ultimately becomes limited in how much it can tax capital. This avoids conflict/tension and therefore supports peaceful democracy
    \item Claims that a demobilized lower class is critical for the peaceful sustenance of authoritarianism
    \item Primary operationalization of capital immobility is land: land cannot be hidden nor moved. Economic modernization involves a shift from immobility (land, natural resources) to mobility (human capital, other mobile capital)
\end{itemize}


\subsection{Important Specific Factual Claims}

\begin{itemize}[leftmargin=*, itemsep=4pt]
    \item
\end{itemize}


\subsection{Connections}

\begin{itemize}[leftmargin=*, itemsep=4pt]
    \item Agrees and expands on Moore
\end{itemize}


\subsection{My Critiques}

\begin{itemize}[leftmargin=*, itemsep=4pt]
    \item
\end{itemize}


% ============================================================
\part{Carey 2009}
\begin{center}
Carey, John M. \textit{Legislative Voting and Accountability}. Cambridge: Cambridge University Press, 2009.
\end{center}

\section{My Notes}

\subsection{Main Idea}



%--------------------------------------------------
\subsection{General Notes}
%--------------------------------------------------

\begin{itemize}[leftmargin=*, itemsep=4pt]
    \item Differentiation between collective and individualistic accountability: an entire legislature or governing party can relatively easily be held accountable by voters for policy \emph{outcomes}, as they are supposedly entrusted with delivering productive results. Yet, it can be difficult for voters to obtain information on specific legislators such that they may be held accountable for their more specific actions. This is a fundamental informational tradeoff delineating two frameworks of legislative accountability.'
    \item Legislators face two primary \textbf{and competing} principals (likely more): their voters can hold them accountable, especially if individualistic accountability is possible, and their party leadership can hold them accountable depending on the institutional strength of the party system. (A high-ranking third principal would be presidents because of their impact on legislative outcomes and processes). These principals may be sequential (i.e., voters $\rightarrow$ parties $\rightarrow$ legislator) or competing (i.e., voters $\rightarrow$ parties $\rightarrow$ legislator AND voters $\rightarrow$ legislator)
    \begin{center}
        \includegraphics[width=13cm]{"Reading Notes/Reading List/Carey 2009/principal models.jpg"}
    \end{center}
    \item Parties become less unified as the number and magnitude of principals increases; unity as an equilibrium outcome instead of a fixed trait
    \item Parties incur costs by disciplining their members -- both from the administration of said costs, but also because it limits individual members' abilities to garner local support which could then more broadly be translated into party support. (This is potentially a Fetterman-esque situation -- if you believe it is more likely that he will vote with the Democrats on some key issues, it may be worthwhile for him to stray further so long as it secures a Democratic seat in Pennsylvania)
    \item Driving question: are votes visible? If individual votes can be tracked and accurately reported, then the likelihood of voters being able to pursue individualistic accountability is much greater (though still contingent on the institutional means of actually doing so -- that is, how they are able to inflict electoral consequence). But, if individual votes cannot be tracked or reported -- for example, via anonymity, voice votes, etc. -- it can become exceedingly difficult for this mechanism to be maintained.
    \begin{itemize}
        \item It is also important to distinguish between \emph{who} can observe these individual votes. In cases where that information may not be made available to voters, it still could be possible for party leaders to identify behaviors internally.
        \item Notes that transparency is not unilaterally a productive thing, as it can weaken collective accountability and (extension) can create adverse incentives against good governance. \textbf{Transparency redistributes accountability instead of plainly increasing it}
    \end{itemize}
    \item Empirical analysis: party unity on roll-call votes, primarily using the RICE index (absolute difference between the shares of a party voting Y/N on a given vote)
\end{itemize}


\subsection{Important Specific Factual Claims}

\begin{itemize}[leftmargin=*, itemsep=4pt]
    \item Legislators face decisiveness problems:
    \begin{itemize}
        \item \underline{Bottlenecks}: There are an infinite (or at least, very large) number of policy options available for any given question, and legislatures must narrow down to some much smaller range of viable alternatives from which to collectively choose. Parties may be a solution to this, in identifying specific policy alternatives as being the general linkage to their own ideological background. However, this also implies limiting the agenda-setting capacity of many legislative actors such that an unbreakable stalemate can be avoided; similarly, veto players must be limited in order to prevent cycling or instability.
        \item \underline{Promotion}: As a legislator becomes more experienced and gains seniority, they both gain access to additional decision-making power and decrease their vulnerability to replacement (this is especially true in party-list elections). This accountability mechanism may subsequently become weaker, requiring alternative forms of punishment/consequence.
        \item ``Party unity is the point of reference because legislators everywhere are partisan, because legislators everywhere answer to party leaders as principals, and because party unity can be identified by the fundamental act of legislative decision making -- that is, voting. Where party unity is lacking, there are various stories we can tell about legislative motivation, but they all involve legislators' making decisions to deviate from the game plan of the team -- the collective unit -- that is the central basis of legislative organization.'
    \end{itemize}
\end{itemize}


\subsection{Connections}

\begin{itemize}
    \item Accountability requires both responsiveness, information about behavior, and the ability to act on that information -- very aligned w/ Powell? (check which author this was)
    \item Critiques typical construction of an SMDP vs. PR tradeoff (explicitly naming Lijphart and Huber/Powell): ``This conclusion rests on two key assumptions, however: that political parties are fundamental units of legislative representation and that a left-right spectrum meaningfully describes the ideological arena of party competition. They are open to greater skepticism in environments where party systems are more volatile or party reputations less stable.'' Claims that the fundamental tradeoff instead should be between whether or not accountability is collective or individualistic (which is itself a downstream impact of the majoritarianism vs. proportionalism debate).
    \item When considering the impact of observability on accountability -- ties very closely to Stokes et al, who suggest that perceiving individual votes (outside of the legislature and instead at the citizenry level) may actually be significantly easier than previously thought, or at least, that voters believe it to be quite possible.
    \item Larreguy \& Raffler 2025 discuss two nested principal-agent problems within developing democracies -- voters and politicians, and politicians and bureaucrats. Contrast however that this is directional rather than doubled-up -- but, the rest of the paper has many citations to be considered here as well
\end{itemize}


\subsection{My Critiques}

\begin{itemize}
    \item 
\end{itemize}


% ============================================================
\part{Cook and Levi 1990}
\begin{center}
Cook, Karen Schweers, and Margaret Levi, eds. \textit{The Limits of Rationality}. Chicago: University of Chicago Press, 1990.
\end{center}

\section{My Notes}

\subsection{Main Idea}



%--------------------------------------------------
\subsection*{Overarching Notes}
%--------------------------------------------------

\begin{itemize}[leftmargin=*, itemsep=4pt]
    \item Overarching: what are the limits of rational choice theory, and what is its future? Claim is that the future of rational choice analysis is based in studies of institutions and norms
    \item Shares the critique/belief that rational choice models are often oversimplified in social science environments where economic principles (price theory, competition, etc.) are not relevant. Extension of game theory to focus on strategic interactions has done some to alleviate this by emphasizing the importance of group decisions, but has not been sufficiently defeated
\end{itemize}


\subsection{Elster: ``When Rationality Fails''}

\begin{itemize}
    \item ``Failure to recognize the indeterminacy of rational choice theory can lead to irrational behavior.''
    \item RCT does not identify what our goals should be, but does identify what should be done in order to achieve some specified goal(s)
    \item Under rationality, beliefs must be optimally related to the evidence available to the agent (kind of quote)
    \item Veyne 1976: ``To serve his master well, a slave must have some independence of execution: beliefs born of passion serve passion badly''
    \item Claim: some goals and utilities are human and cannot be negated even in the search for greater utility (for example, the desire for freedom cannot be broadly eliminated)
    \item Hyperrationality: ``the failure to recognize the failure of rational choice theory to yield unique prescriptions or predictions.''
    \item ``Rationality is presupposed by any competing theory of motivation, whereas rationality itself does not presuppose anything else. On grounds of parsimony, therefore, we should begin by assuming nothing but rationality.''
    \item Social norms can be seemingly irrational, but it is perhaps following them which produces rationality -- norms may be the strongest alternative to RCT if you assume that it does not itself result in rationality
\end{itemize}


\subsection{Tversky \& Kahneman: ``Rational Choice and the Framing of Decisions"}

\begin{itemize}
    \item Deviations from predicted rational outcomes are significant, potentially disqualifyingly so; concludes that the normative and descriptive analyses ``cannot be reconciled''
    \item Von Neumann Morgenstern utility: cancellation, transitivity, dominance, and invariance
\end{itemize}


\subsection{Machina: ``Choice Under Uncertainty: Problems Solved \& Unsolved''}

\begin{itemize}
    \item Questions arising about theories of choice under uncertainty, which has come to face a handful of challenges
\end{itemize}


\subsection{Stigler \& Becker: ``De Gustibus Non Est Disputandum''}

\begin{itemize}
    \item Foundational rooting in tastes as an exogenous and generally-immovable characteristic
    \item ``We are proposing the hypothesis that widespread and/or persistent human behavior can be explained by a generalized calculus of utility-maximizing behavior, without introducing the qualification `tastes remaining the same.'''
    \item Rational choice as the best available tool to explain behavior without making assumptions that tastes have changed -- that is, that allows for (but does not require) consistent tastes
    \item Fashion as a counterexample; ``The obvious method of reconciling fashion with our thesis is to resort again to the now familiar argument that people consume commodities and only indirectly do they consume market goods, so fashions in market goods are compatible with stability in the utility function of commodities.''
    \item General argument is that you can reconcile changing behavior with static tastes
\end{itemize}


\subsection{Taylor: ``Cooperation and Rationality: Notes on the Collective Action Problem and Its Solutions''}


\begin{itemize}
    \item Elster's definition of collective action problem:
    \begin{itemize}
        \item Strong: Prisoners' Dilemma
        \item Weak: (1) everyone prefers universal cooperation to universal noncooperation, and (2) cooperation is ``individually unstable and individually inaccessible.''
    \end{itemize}
    \item Two types of collective action solutions:
    \begin{enumerate}
        \item Spontaneous / Internal: no changes to the game, preferences, or beliefs
        \item External: changing the game, preferences, or beliefs
    \end{enumerate}
    \item External solutions can be either decentralized, such that the initiative or responsibility for enacting the solution is broadly shared, or centralized, such that the solution authority is concentrated in the hands of some subset of the members (assume some sort of domination outcome here?). Centralized is taken to be consubstantial with the state
    \item Highlights the importance of repetition, group norms, and an ability to enforce said norms as a core component of solving collective action problems
    \item One external solution is the political entrepreneur: coordination mechanism. Must change beliefs or preferences (usually occurs by adding knowledge)
    \item Property rights as a possible solution to collective action problems: if the impacts of consumption can be totally internalized within the individual, then such problems may be avoided entirely
    \item All seems very similar to Ostrom's operationalization of the Folk theorem
\end{itemize}


\subsection{Coleman: ``Norm-Generating Structures''}


\begin{itemize}
    \item Provides a comprehensive (but rather esoteric) definition of norms
    \item ``The demand for a norm will arise when an action by one actor imposes externalities on other actors. When the externalities are positive, the demand is for a prescriptive norm, to encourage or induce the focal action. When the externalities are negative, the demand is for a proscriptive norm, to discourage the focal action.''
    \item Norms emerge as behaviors within (and surrounding) repeated PD interactions. Question then becomes: how can these be supported when repeated interactions are not guaranteed? Requires some formation of a partial tendency to not defect in order to support broader social cooperation -- i.e., some initial impetus that can then be built upon
\end{itemize}


\subsection{Stinchcombe: ``Reason and Rationality''}

\begin{itemize}
    \item Specifies an ideal long-run relationship between Weber's substantive and formal rationality:
    \begin{itemize}
        \item Formal rationality: ``standardized methods of calculation on which routines can be based'' -- strict and explicit institutions
        \item Substantive rationality: Beyond formal methods and instead engaging with substance -- how strong is your CV, is monetary gain a sufficient incentive, etc.
    \end{itemize}
    \item ``What we want of our institutions of reason is a proper balance between efficient formal approximations that can have a reliable social effect, and substantive good sense to know their limits and to improve them.''
\end{itemize}


\subsection{Miller: ``Managerial Dilemmas: Political Leadership in Hierarchies''}

\begin{itemize}
    \item Two distinct literatures on organizational control:
    \begin{enumerate}
        \item Control by means of designing aligned incentive systems and sanctions such that self-interested / utility-maximizing actors behave in some desired way
        \item Control by means of behavioral influence -- inspiration, social encouragement, generally means that are not tied to some form of explicit payoff or sanction
    \end{enumerate}
    \item General claim: this first approach is self-destructive, and is demonstrated by social choice theory to be logically impossible. The most effective approach to control is subsequently personalistic leadership (tie to Weber)
    \item ``Once the Leviathan is created to correct the deficiencies of collective action, Leviathan itself is found to generate inefficiencies. The very mechanism that makes Leviathan desirable -- the authority to change members' incentives -- is found to be logically insufficient to guarantee efficiency. The irony, then, is that the managers of a hierarchy will find it desirable to look for modes of inducing voluntary cooperation within the Leviathan.'
\end{itemize}


\subsection{Hardin: The Social Evolution of Cooperation}

\begin{itemize}
    \item Coordination as a foundational principle for cooperation -- we must be able to coordinate on some cooperative process, if we are to achieve (and sustain) that cooperation
    \item Gunman theory of law: ``if we are to be made to obey law by threat of force, then the state will be unable to mount adequate mechanisms of enforcement.''
    \item Perverse logic of state construction: if we construct the state to enforce collective action solutions, then we must overcome the collective action problem of constructing the state itself
    \item Common law as a reflection of path-dependency
    \item Coordination creates a convention, and a convention promotes further coordination -- circular pathway
\end{itemize}


\subsection{North: ``Institutions and Their Consequences for Economic Performance''}

\begin{itemize}
    \item Largely consistent with North 1990a -- institutions reduce transaction costs
    \item ``Institutions consist of informal constraints and formal rules and of their enforcement characteristics.''
    \item Effectiveness of institutions is largely reflected through how much they reduce transaction costs
    \item Cites North 1981 -- two reasons for institutions producing lesser property rights. First, may pander to certain powerful groups who may otherwise support an alternative ruler/regime. Second, may seek to gain tax revenue (which could be decreased through stronger property rights)
    \item North 1990b: ``a basic function of institutions is to provide stability and continuity by dampening the effects of relative price changes. It is institutional stability that makes possible complex exchange across space and time.''
\end{itemize}


\subsection{Levi: ``A Logic of Institutional Change''}


% ============================================================
\part{Cox 1987}
\begin{center}
Cox, Gary W. \textit{The Efficient Secret: The Cabinet and the Development of Political Parties in Victorian England}. Cambridge: Cambridge University Press, 1987.
\end{center}

\section{My Notes}

\subsection{Main Idea}



%--------------------------------------------------
\subsection{General Notes}
%--------------------------------------------------

\begin{itemize}[leftmargin=*, itemsep=4pt]
    \item ``We seek to explain the evolution of two crucial behavioral aspects of the modern English polity: the marked regularity with which MPs vote with their parties, and the strong tendency of voters to vote for parties rather than individuals.... To show when and explain why party became the dominant influence on voting behavior in England -- both in the legislature and in the electorate.''
    \item Two primary independent variables: ``the extension of the suffrage and the centralization of legislative authority.''
    \item Rapid expansion of private bills (see below section) as the 19th century progressed, largely stemming initially from the complexity of regional investments in railroad services. A great deal of lobbying became the norm as voters sought to win the favor of MPs, but as this became broadly known, also led to significantly more public debate through procedural means
    \item Mid-century parliament (post first and second Reform Acts) featured a great deal of party independence: legislators stood to gain significant clout by voting their conscience rather than their party, and party discipline plummeted. Yet, it subsequently significant rebounded later in the 19th century
    \item Primary hypotheses:
    \begin{enumerate}
        \item ``Since the cost of becoming informed increases with the complexity and frequency of votes, nineteenth century MPs should have been more likely to rely on free information -- a chief element of which was the recommendations of their party leaders -- as motions pushed to a division became more complex and frequent."
        \item ``Since the benefit to acquiring information depends on the chances that the information will actually change the vote decision, any process whereby nineteenth century MPs became more confident that their party's line would be acceptable to them were they to investigate the matter thoroughly should have increased party voting. Party cohesion should have increased when the probability of agreement between the party and other relevant groups increased.''
    \end{enumerate}
    \item Argument: as the information available to voters greatly increased, and the number of voters for whom that information was relevant increased as well, it became an informational challenge for voters to hold specific MPs accountable; instead, voters turned towards the parties as broader symbols of accountability, to be punished and rewarded collectively. Following this transition, it became incentive compatible for individual MPs to cede power/authority to the collective party, as its success would in turn secure their own position(s).
\end{itemize}


\subsection{Important Specific Factual Claims}

\begin{itemize}[leftmargin=*, itemsep=4pt]
    \item Private bill: dealing with a specific policy issue regarding a firm, individual, constituency, etc. -- NOT designed to establish blanket policy guidelines, but rather to solve a particular problem within the MP's constituency (explicitly or implied)
    \item Core developments in 19th-century Britain:
    \begin{itemize}
        \item Reform Act of 1832 \& Reform Act of 1867 provided significant expansions of the franchise and provided greater equality of that franchise (previously -- and still following 1832 -- many small constituencies were heavily overrepresented in Parliament). Further expansion and ``clean-ups'' in the third Reform Act
        \item Dramatic increase in population met by an unexpected proportional increase in real wealth: industrial revolution allowed for greatly increased productivity and output per worker and in turn lowered consumer costs \& facilitated broader consumer markets. Shift in population from agrarian to industrial focus
        \item Rapid expansion of number of news outlets and the number of citizens exposed to the press: lowered/eliminated duties on advertisements/stamps/papers, steam-driven printing press in 1811, development of the telegraph in the late 1840s. Rapid diffusion of political information, particularly during elections
        \item Development/expansion of steam locomotive and developed railroad system, providing both significant economic growth (\& complexity) and interpersonal transportation; this in turn forced procedural reforms in parliament with the a delineation between ``public'' and ``private'' bills (think: who benefits?).
    \end{itemize}
    \item The ``efficient secret'' itself is the fusion of legislative and executive authority within the cabinet, rather than as separate authoritative bodies. Core question is then what facilitated this delegation of power from individual MPs to their parties and, subsequently, the respective cabinet.
\end{itemize}


\subsection{Connections}

\begin{itemize}
    \item Olson 1965: party formation as a collective action solution. Each MP could retain their own influence/power/oversight, but in doing so would greatly hinder the collective legislative process -- individuals needed to be incentivized to acquiesce these individual benefits for the good of the party as a whole. (Is this directly cited?)
    \item Huntington 1968: political development as a response to economic/social development, very cleanly reflected in Cox's account of parliament
    \item Crisp et al 2004: legislators deciding to focus on either local or national issues -- either providers of public goods or providers of private goods but to their constituents. Connected to the ``private bill'' whereby legislators can choose to focus on more specifically/local policy concerns
    \item Carey \& Shugart 1995 -- under which electoral systems would legislators be more likely to pursue a personal vote (and subsequently less likely to cede power to the party / cabinet)?
    \item Jones \& Hwang 2005, Katz \& Mair 1995 -- Cartel theory: Cartel theory claims that the party has very limited control over its members; Cox suggests that this has been largely overcome. A driving question here (and good to write about) -- is this because cartel theory is disproven, or is it because it is still incentive-compatible for MPs to cede authority to the central party/cabinet, such that cartel theory is circumvented in a new equilibrium outcome?
\end{itemize}


\subsection{My Critiques}

\begin{itemize}
    \item
\end{itemize}


% ============================================================
\part{Cox 1997}
\begin{center}
Cox, Gary W. \textit{Making Votes Count: Strategic Coordination in the World's Electoral Systems}. Cambridge: Cambridge University Press, 1997.
\end{center}

\section{My Notes}

\subsection{Main Idea}



%--------------------------------------------------
\subsection{General Notes}
%--------------------------------------------------

\begin{itemize}[leftmargin=*, itemsep=4pt]
    \item Contributions: Elections as a coordination problem, empirical tests based on constituency (rather than national aggregate) level.
    \item Elections represent coordination problems: parties (and voters) must coalesce around certain candidates, and do so both by and in light of strategic voting. Yet, this strategic coordination does not always have clear outcomes for winners/losers. Successful coordination often requires an ideological tradeoff that is often unappealing to strictly-ideological factions.
    \item Elections are by and large market affairs, where ``market-clearing'' expectations balance the supply and demand for a number of candidates. Yet, election behavior does not always follow equilibrium outcomes, and in fact a great deal of electoral politics happens outside of these confines or in otherwise unstable institutional settings
    \item Duvergian critiques: broad sociological critiques largely stemming from accusations of oversimplification or ``institutional determinism'': institutions determine this outcome regardless of cultural or social norms, behavior, context, etc.
    \begin{itemize}
        \item Is the direction of causality correct? Could it be the case that party systems determine electoral systems, rather than the inverse? Citations listed on pg. 15
        \item Is the electoral system truly the variable of interest, or might more emphasis fall on the number and structure of social cleavages?
    \end{itemize}
    Quite a bit of this also rises out of questions of endogeneity over electoral systems -- they are, of course, endogenous in initial construction, but what about their long-term forms? Presumably parties can shape electoral institutions over time in accordance with their own strategic visions
    \item Three stages when accounting for the level of vote or seat concentration in some polity:
    \begin{enumerate}
        \item Translation of social cleavages into partisan preferences;
        \item Translation of partisan preferences into votes;
        \item Translation of votes into seats.
    \end{enumerate}
    \item Core construction is the ``$M + 1$ Rule'': if a district offers $M$ seats, we should expect to see consolidation down to $M+1$ viable candidates/parties/etc. participating in that election. Duverger's Law is subsequently a simplification of this stratified between $M = 1, M > 1$ (majoritarianism vs. proportionalism, respectively).
\end{itemize}


\subsection{Important Specific Factual Claims}

\begin{itemize}[leftmargin=*, itemsep=4pt]
    \item Identifies an \emph{a priori} expectation on the number of parties competing, compared between legislative chambers within a country
    \item Limitation of Duverger: what if cleavages are geographically concentrated? A party can have a stronghold within a small region, but still be relegated to third-party status overall. Think SNP; Duverger's should predict a combination with Labour (but is this potentially stymmied by the possibility of coalition formation?)
\end{itemize}


\subsection{Connections}

\begin{itemize}
    \item Heavy framing based on the Gibbard-Satterthwaite theorem, but notes that its result -- though quite rigorous -- can be a bit more difficult to apply to political science (just provides a binary Y/N on strategic voting)
    \item Highlights research gap in interaction between electoral systems and social cleavages (though may have been filled in past 30 years)
\end{itemize}


\subsection{My Critiques}

\begin{itemize}
    \item
\end{itemize}


% ============================================================
\part{Dahl 1972}
\begin{center}
Dahl, Robert A. \textit{Polyarchy: Participation and Opposition}. New Haven: Yale University Press, 1971.
\end{center}

\section{My Notes}

\subsection{Main Idea}



%--------------------------------------------------
\subsection{General Notes}
%--------------------------------------------------

\begin{itemize}[leftmargin=*, itemsep=4pt]
    \item ``Given a regime in which the opponents of the government cannot openly and legally organize into political parties in order to oppose the government in free and fair elections, what conditions favor or impede a transformation into a regime in which they can? That is the question with which this book is concerned.''
    \item Democratization and development of a public opposition are separate forces. We consider two metrics: the extent to which an opposition is allowed to reasonably exist, and the size of the franchise granted government participation. Two dimensions become public contestation and the right to participate
     \begin{center}
        \includegraphics[width=10cm]{"Reading Notes/Reading List/Dahl 1972/generalframework.jpg"} \\
        Democratization represented by path III
    \end{center}
    \item ``The greater the conflict between government and opposition, the more likely that each will seek to deny opportunities to the other to participate effectively in policy making. To put it another way, the greater the conflict between a government and its opponents, the more costly it is for each to tolerate each other''
    \begin{center}
        \includegraphics[width=10cm]{"Reading Notes/Reading List/Dahl 1972/dahl 1.4.jpg"}
    \end{center}
    \item As the franchise expands, candidates change their alignments/platforms/beliefs: need to seek support of the newly-enfranchised groups. May often result in the formation of new parties (especially workers'/Labour parties or the like, as well as socialist sentiment) to appeal to these groups. Also drives a change in the structures of parties, as they need to be able to better connect with a larger group: introduces hierarchy, bureaucracy, etc.
    \item Three primary paths to polyarchy:
    \begin{enumerate}
        \item \underline{Liberalization precedes inclusiveness}: England, Sweden. Most common, more peaceful as groups slowly integrate and existing ties hold cohesion. Supported by mutual security that exists at the beginning, and then can be maintained. (Very close to A\&R ``Why the West...'')
        \item \underline{Inclusiveness precedes liberalization}: Wiemar Republic. Less common and more dangerous: mutual security system does not exist at the beginning of the process, and so instability may perpetuate.
        \item \underline{Shortcut}: closed hegemony abruptly transforms into a polyarchy by a sudden grant of suffrage and rights of contestation. France 1789-1792. Few examples of working; no mutual security and little time to adapt to things. (Extremely reflective of Huntington argument, although Dahl focuses more on the lack of mutual security rather than a lack of expressive participatory opportunities)
    \end{enumerate}
    \item Polyarchies often ``inaugurated'' either within an independent nation-state (transformed via reform, revolution, or military conquest) or in a dependent country subject to another state (transformed via reform or revolution). Peace lends future stability: polyarchies developed through reform are more likely to be viewed as legitimate
    \item Because suppression requires resources (both for violent and nonviolent means), regimes with limited resources may be more likely to tolerate the opposition, as the relative costliness of intolerance increases. This is counterintuitive in many senses to existing arguments that wealth would drive stability: rather, wealth drives state capacity, which in turn allows for repression/hegemony.
    \item Effective and dispersed government requires that means of repression/violence/etc. are available to \textit{both} the government and the opposition: asymmetric availability results in monopolistic governance while symmetric unavailability results in uniformity (but could perhaps be more unstable, depending). Ideal for dispersement/neutralization for \textit{both} access to violence and access to socioeconomic sanctions; hegemony thrives when both are monopolized
    \item While classic liberal theory suggested that competitive politics could not be sustained without competitive market-based economics, Dahl disagrees that this universally holds true -- rather, the motivating factor in hegemony is whether or not there is a highly centralized economy (which is itself dissuaded by competitive politics)
    \item ``The chances that a country will develop and maintain a competitive political regime (and, even more so, a polyarchy) depend up on the extent to which the country's society and economy:
    \begin{enumerate}
        \item Provide literacy, education, and communication
        \item Create a pluralistic rather than a centrally dominated social order
        \item Prevent extreme inequalities
    \end{enumerate}
    among the politically relevant strata of the country."
    \item Polyarchy arises out of advanced economies, but advanced economies also may be more prone to severe inequality.
     \begin{itemize}
        \item One argument here invlves the distribution of political resources and skills. Those with the means to do so will attempt to bend political systems in their favor; the distribution of those means is itself endogenous and largely favors to those with the means in the first place. Advanced economies lead to less inequality of political resources -- does not imply strict equality, but does imply parity, which in turn supports both inclusion and contestation
        \item Alternative explanation suggests that the driving force here is the power of resentment from excluded groups. Regimes may capitulate to \emph{some} demands for greater equality, but not necessarily all; doing so may be sufficient to gain the support of excluded groups even if full equality is not actually achieved. Possible ``performative'' or signaling impact here
    \end{itemize}
    \item The effectiveness of polyarchy is largely driven by a belief in its legitimacy. Three propositions here:
    \begin{enumerate}
        \item ``The greater the belief within a given country in the legitimacy of the institutions of polyarchy, the greater the chances for polyarchy.''
        \item ``Countries vary a great deal in the extent to which activists (and others) believe in the legitimacy of polyarchy.''
        \item ``These variations are to some extend independent of the social and economic characteristics of a country: two countries with a great many similarities in their social and economic orders may vary significantly in the extent to which activists (and others) believe in the legitimacy of polyarchy.''
    \end{enumerate}
    \item Trust is a fundamental component of a successful polyarchy. Citizens must be able to effectively communicate, which requires trust. But, more notably, polyarchy requires that political actors trust their opponents (essentially lowering the stakes of political loss): a trust that if the opposition takes power, they will neither impede my ability to lawfully gain back power in the future nor violate my rights/standing/etc. in the meantime. The stakes \emph{need} to be lowered in order to sustain cooperation.
\end{itemize}


\subsection{Important Specific Factual Claims}

\begin{itemize}[leftmargin=*, itemsep=4pt]
    \item Three necessary conditions for democracy -- citizens must be able to:
    \begin{enumerate}
        \item Formulate their preferences
        \item Signify their preferences to their fellow citizens and the government by indvidual and collective action
        \item Have their preferences weighed equally in the conduct of the government -- weighted with no discrimination because of the content or source of the preferences
    \end{enumerate}
    \item Regime type has consequences. (Argues specifically for this claim in chapter 2 -- although quite obvious, may be a useful early example of such)
    \item When military capacity is high and organized, polyarchy is impossible unless the military is sufficiently depoliticized such that it has neither intent nor ability to interfere with political discourse or dynamics.
    \item Farming societies largely differentiated into traditional peasant society, defined largely by hegemony, or a free farmer society, more aligned towards polyarchy.
    \item Though socioeconomic development seems in part to be linked to the prevalence of polyarchy -- and it seems plausible that there exist upper and lower thresholds of economic development beyond which the probability of polyarchy remains constant -- the relationship is neither clear, nor statedly causal, nor incredibly strong, nor without exception.
    \item Developed economies allow workers to be literate and more-educated, which in turn can expand their potential inclusion into the polyarchy -- dual-directional effect. Advanced economies cause pluralistic social orders, which in turn cause demands for a competitive political system (under one causal chain).
    \item ``In an advanced economy, long-run performance under threat or coercion is less productive at all levels than a more willing performance based upon voluntary compliance.''
    \item Polyarchy seems to rest largely on mixed cooperative-competitive relationships: cooperate on the aspects of policy and governance whereupon common ground can be found, and compete on the others. In the absence of some cooperation, the trust and cohesion required for polyarchy cannot be sustained. But, in the absence of some conflict, the ability to retain a competitive political dynamic is greatly reduced.
    \item The historical association of polyarchy with capitalism called its value into question following the Great Depression, as its legitimacy was partially derailed. The explosion of single-party systems in the wake of this -- including the acceleration of previously unstable single-party systems -- can be partially attributed to this, as previous claims that polyarchy and democracy supported economic development became less clear.
    \item Three primary impacts of foreign involvement on polyarchy:
    \begin{enumerate}
        \item Foreign actors can affect the beliefs of domesti actors, which we have already shown to be critical in establishing the conditions for (or against) polyarchy
        \item Foreign actors can affect the decision sets available to domestic political leaders; for example, Cold-War-style tension between US and USSR affecting the political outcomes of dependent countries
        \item Foreign actors can outright impose regime types or political systems on a domestic nation, either through conquest, coercion, or other similar means of domination
    \end{enumerate}
\end{itemize}


\subsection{Connections}

\begin{itemize}
    \item ``When hegemonic regimes and competitive oligarchies move toward polyarchy they increase the opportunities for effective participation and contestation and hence the number of individuals, groups, and interests whose preferences have to be considered in policymaking.'' Indirect connection to Huntington 1968: increases participation, which per Huntington should result in instability unless there are additional institutions to handle it. Suggestion is that transition to polyarchy must involve a non-trivial movement towards institutional development as well.
    \item More Huntington 1968: agrees (indirectly) that parties begin to move away from tribal/familial/otherwise native sources of power and authority, and instead are replaced by those which are appealing to the middle classes
    \item Broad-scale oppression (i.e., Holocaust, USSR atrocities, etc.) do not happen in polyarchies, largely because the low barriers to contestation both allow marginalize groups to better defend themselves politically and preemptively dissuade the state from pursuing or proposing any such aggressive policies
    \item Liberalization preceding inclusiveness -- very much tied to A\&R Why the West Extended the Franchise; essentially the same mechanism. Shortcut -- very much tied to Huntginton, although Dahl focuses more on mutual security than effective participation. See ``three paths to polyarchy'' in second section.
    \item Discussion of polyarchy as potentially losing legitimacy in the wake of the Great Depression, as previously it seemed to be a given that democracy drove economic development in some unilateral way -- certainly goes back to Huntington, and also provides potentially some empirical backing to Huntington's argument that democracy alone did not drive economic development (as well as the notion that voters actually internalized that lapse).
    \item Four primary factors in whether or not an individual adopts new political beliefs:
    \begin{enumerate}
        \item Exposure: how much are they actually exposed to these new ideas? Is their political community well connected and civically engaged enough that this ``contagion'' could actually take place?
        \item Prestige: how prestigious is the new idea? Who are its proponents, and how much credit is given to their endorsement? Have the ideas been well-supported and either proven or disproven in high-profile settings?
        \item Consistency w/ beliefs: are the new beliefs consistent with the individual's present beliefs? Tautologically, there must be some separation for this to be a meaningful shift, but are they aligned in principle or entirely separated? Do they fit into a common worldview and perhaps merely offer a differential explanation? (RG note: this could largely be tied to religious impacts on political beliefs, most notably Catholicism in modern western politics)
        \item Consistency w/ experience: are the new beliefs consistent with the policy-related (or even simply emotional) experiences the subject has endured? Do they align with previous needs/wants/solutions
    \end{enumerate}
    \item See Boix 1999: ruling parties can use the transition between majoritarianism and PR as a way of maintaining some degree of exclusivity (but possibly increasing contestation in part)
    \item Gourevitch 1978 expands particularly on the influence of foreign politics on domestic structures -- Dahl notes that invasion can be a key factor in domestic politics, while Gourevitch acknowledges and significant expands upon this
    \item Schmitter 1974 discusses corporatism as a form of hegemonic politics and a limited range of inclusiveness
\end{itemize}


\subsection{My Critiques}

\begin{itemize}
    \item One central claim seems to be that competitive politics universally dissuades coercion/repression etc. There is a world, though, where that argument does not necessarily stand. For example, in a reverse causal chain, if coercive capacity was quite high -- such that the rents of holding office were similarly high -- we might expect there to be more competition endogenously.
\end{itemize}


% ============================================================
\part{Huber and Shipan 2002}
\begin{center}
Huber, John D., and Charles R. Shipan. \textit{Deliberate Discretion? The Institutional Foundations of Bureaucratic Autonomy}. Cambridge: Cambridge University Press, 2002.
\end{center}

\section{My Notes}

\subsection{Main Idea}



%--------------------------------------------------
\subsection{General Notes}
%--------------------------------------------------

\begin{itemize}[leftmargin=*, itemsep=4pt]
    \item Though voters care about policy outcomes, they are generally agnostic towards the actual writing of the laws that drive them. These statutory details are core drivers of policy outcomes, but are therefore largely ignored by the general public. As a result, politicians have two general strategies in achieving policy ends through statutory means:
    \begin{enumerate}
        \item ``Write long statutes with extremely detailed language in an effort to micromanage the policymaking process''
        \item ``Write vague statutes that leave many details unspecified, thereby delegating policymaking authority to other actors, usually bureaucrats''
    \end{enumerate}
    \item ``Our argument draws attention to four general factors that affect politician's incentives and ability to use legislation to micromanage policy-making:"
    \begin{enumerate}
        \item The level of policy conflict between the politicians who adopt the statutes and the agents who implement them;
        \item The capacity of politicians to write detailed statutes;
        \item The bargaining environment in which statutes are adopted;
        \item the expectations that politicians have about nonstatutory factors (course or legislative oversight opportunities), independent of the language of the statutes
    \end{enumerate}
    \item Good governance essentially requires the construction of the \emph{administrative state}, but the administrative state itself is a somewhat undemocratic phenomenon: accountability is indirect, and power is vested in hands that are not necessarily elected or even closely impacted by elections. This is particularly notable given the high level of power/authority vested within bureaucratic systems, who are responsible for the vast majority of public policymaking in modern societies.
    \item Frame statutes as ``blueprints'' of public policy: set out the framework and guidelines that bureaucrats will subsequently follow, and in doing so choose a degree of specificity / vagueness to in turn yield some degree of autonomy.
    \item To a large extent, statutory specificity depends on \emph{legislative capacity}, which in turn ``affects the costs to politicians of drafting detailed legislation.'' Costs are information and opportunity-related (time, other resources, etc.) Highly varied at state legislature levels.
    \item ``If politicians can depend on other nonstatutory features of the political environment to help produce bureaucratic actions that are congruent with their preferences, this diminishes the extent to which precise statutory language is necessary for obtaining desired policy outcomes.''
    \item Two core mistakes possible from this arrangement:
    \begin{enumerate}
        \item \underline{Bureaucrat given too much discretion}: Bureaucrat is empowered to choose a policy yielding an outcome that matches their own ideal point, but may be far from the politician's ideal point
        \item \underline{Bureaucrat given not enough discretion}: Because bureaucrat has additional (and private) information about policy-outcome mapping, they may know that delivering the politician's ideal outcome requires enacting a policy far from the politician's ideal policy -- yet, if they are not sufficiently empowered to stray that far, they may be bounded in how close they can deliver to the ideal outcome
    \end{enumerate}
    \item Arguments as to how the level of discretion in legislation should vary with different politican and institutional contexts (all direct quotes):
    \begin{itemize}
        \item As policy conflict between politicians and bureaucrats increases, or as legislative capacity increases, the relative value to politicians of low-discretion statutes increases.
        \item Whether the model predicts a direct or interactive effect of these two variables on discretion depends on the distribution of legislative capacity and policy conflict among the political systems that are analyzed. If there are a large number of political systems that never have sufficient capacity to adopt a low-discretion law, then we should expect interactive effects (where, for example, legislative capacity influences the design of statutes only if the level of policy conflict is sufficiently high). By contrast, if a large numbere of political systems have high legislative capacity, the effects of these two variables should be direct.
        \item As the ability of politicians to rely on nonstatutory factors to achieve policy objectives increases, the relative value of low-discretion laws declines.
        \item These results hold in parliamentary and presidential systems.
        \item In presidential systems, bicameral conflict often leads to more discretion.
    \end{itemize}
    \item ``When conflict is absent and legislators trust policymakers in the agency, they are likely to design vague statutes that allow the agency to draw on its expertise to make policy. And when conflict is present, legislators who disagree with the agency's policy views will constrain the agency by putting specific wording into legislation.''
    \item ``When a different party controls each chamber, the party favoring the governor has the incentive and ability to block some -- but not all -- legislative language that ties the hands of the relative agency.''
\end{itemize}


\subsection{Important Specific Factual Claims}

\begin{itemize}[leftmargin=*, itemsep=4pt]
    \item ``Wise and salutary neglect'': Edmund Burke's 1775 \emph{Speech on the Conciliation of America} -- providing a broad latitude to operate independently and subsequently create individual laws, procedures, etc. Synonymous with high political autonomy
    \item A core factor in the necessity of the administrative state is an informational problem: we need technical expertise for certain tasks that cannot necessarily be centralized in the elected political system. This requires providing some large degree of autonomy, which in turn invites risk (think principal-agent problem)
    \item Model largely built on informational asymmetry: there is some mapping from policy to outcome, which is perfectly known by the bureaucrat, but is not known by the politician. Both care about outcome, but politician can therefore not perfectly match it, and requires the bureaucrat's information.
    \item Prime case study is Michigan's Medicaid expansion in the late '90s: \\
    \begin{quote}
        There was, however, an alternative to having the state set prices itself. The state realized that although it did not know the marginal costs and benefits faced by the providers, someone else did -- namely, the providers themselves. These providers knew better than anyone else what the true costs were. The question then became how to get this information from the providers. Furthermore, in order for these providers to know the costs and benefits of participation, they would have to know exactly what services they would need to provide. And this, needless to say, would require addressing all the technical issues just discussed.
    \end{quote}
    \item Another quote on single-party dominance: \\
    \begin{quotation}
        The legislator and Havemann were having a meal together in April and she said, ``We should get together and talk.'' And his reply was, ``I don't know why we need to talk to you. You have no way to stop us if we want to do something.''
    \end{quotation} 
    \item On Michigan case: ``Two features of this change in the design of statutes are worth noting. First, the constraints that appeared in the statutes were largely narrow descriptions of policy rather than procedural constraints on agency decision-making processes. Procedures were clearly present and non-trivial, but the participants we talked to emphasized the policy restrictions, which represent the lion's share of the new legislative language. Second, although the number of policy constraints rose when the Democrats took over the House, the Democrats were not able to insert all the policy constraints they desired into the statutes. Both our interviews and our examiniation of the initial House version of the MMC legislation in 1997 make it clear that several of the Democrats' desired policies had to be left out of the final version of the bill."
\end{itemize}


\subsection{Connections}

\begin{itemize}
    \item Where previous literature in autonomy generally focuses on a single case at a time, Huber \& Shipman provide a novel comparative perspective so that the institutional environment can be evaluated as a variable (rather than a fixed component of the situational context).
\end{itemize}


\subsection{My Critiques}

\begin{itemize}
    \item
\end{itemize}


% ============================================================
\part{Huntington 1968}
\begin{center}
Huntington, Samuel P. \textit{Political Order in Changing Societies}. New Haven: Yale University Press, 1968.
\end{center}

\section{My Notes}

\subsection{Main Idea}

When economies and societies modernize, the resulting increase in political conciousness can lead to unrest and instability as participation outpaces institutionalization. To maintain order, political institutions must modernize as well, developing sufficiently as to satisfy this increased participation. This goes against previous intuition (the liberal or ``Whig'' argument) suggesting that economic growth should organically result in political stability, and suggests that political institutionalization must be intentional and structured. Institutional development may fall into revolution when existing state capacity is low and calls for change are met without an institution to consider and implement them. Conversely, well-developed and strong institutional structures may be able to change via reform instead, whereby development follows from the existing institutional structures and occurs in a more controlled manner.


%--------------------------------------------------
\subsection{General Notes}
%--------------------------------------------------

\begin{itemize}[leftmargin=*, itemsep=4pt]
    \item Social and economic development preclude political development: just having economic growth does not improve political outcomes unless institutions are also formed to shape item
    \item When social structures are strong or homogenous, they may take the form of political institutions -- but, in their absence (or when there are many), we need political institutions to effectively coerce government
    \item Madison Federalist 51: ``In framing a government which is to be administered by men over men, the great difficulty lies in this: you must first enable the government to control the governed; and in the next place oblige it to control itself.''
    \item Mature institutions have their own goals, incentives, traditions, etc. -- they no longer simply serve those who created them, but they further serve themselves
    \item Argues that what is good for institutions is usually good for the country and is in some way a representation of the ``public good'' -- that institutional strength is commensurate with effective governance
     \begin{itemize}
        \item Subclaim: governments are ``good'' when their actions reflect the desires/goals of their institutions. \emph{Alignment between institutions and their leaders/operators }
    \end{itemize}
    \item Trust inhibits organization, which inhibits political stability
    \item Modernization involves an openness to new ideas and change, and an expansion of knowledge about one's environment:
    \begin{enumerate}
        \item  ``Political modernization involves the rationalization of authority, the replacement of a large number of traditional, religious, familial, and ethnic political authorities by a single secular, national political authority.'' Man begins to rationalize and question authority, requiring power to be centralized with reason rather than by tradition
        \item ``Political modernization involves the differentiation of new political functions and the development of specialized structures to perform those functions.'' Complexity is tied to modernization; recognize the need for certain tasks
        \item ``Political modernization involves increased participation in politics by social groups throughout society.'' Political participation broadens
    \end{enumerate}
    \item Unrest is tied towards the speed of modernization: rapid development lends instability while graduate development can allow for a moderated and more peaceful transition
    \item Political unrest \emph{requires} economic growth: in total abject poverty, unrest essentially cannot exist (think oxygen-starved)
    \item Three-step relationship:
    \begin{enumerate}
        \item Social mobilation outpacing economic development $\rightarrow$ social frustration
        \item Social frustration outpacing mobility opportunities $\rightarrow$ political participation
        \item Political participation outpacing political institutionalization $\rightarrow$ political instability
    \end{enumerate}
    \item Corruption occurs in response to this last tradeoff: political participation outpaces opportunities for participation, and so participation becomes unstructured and disorganized
    \item Modernization breeds corruption in three ways: core point is that not only do behaviors change, but perceptions of those behaviors do as well. Moving target of what constitutes corruption
    \begin{enumerate}
        \item Modernization involves a change in the basic values of society -- new behaviors may have previously been seen as corrupt, or vice versa. Changing expectations for the role of public figures and institutions. Also ties back to earlier points about questioning legitimacy of authority. May be a result of blurred line between public and private interests
        \item Modernization creates new sources of wealth and power whose relation to politics is not presently defined
        \item ``Changes it produces on the output side of the political system'' -- generally leads to more laws, which themselves offer more opportunity for corruption at every stage
    \end{enumerate}
    \item Top-heavy corruption is commensurate with very low political development, but developed political systems often have bottom-heavy corruption (i.e., corruption is much more severe with local officials and institutions than national). US an example of the former; Trump era pushing towards the latter
    \item Corruption can speed economic development: brings in groups that would otherwise be marginalized, push through bureaucratic oversight, etc. US in 1880s-1890s is key example (think large monopolies)
    \item Defines a ``Green Uprising'' process. First, urbanites develop and unseat rural elites. Then, violence/tension forms between ruling urbanites and remaining rural population en masse. Green Uprising refers to the alliance of some sector of the rural population with some sector of the urban population -- i.e., a workers' movement -- that subsequently overcomes these impasses and brings stability. However, forms of this uprising may differ, particularly in (a) who is part of the movement and (b) what the goals/targets of the movement may be
    \item ``A society is vulnerable to revolution only when the opposition of the middle class to the political system coincides with the opposition of the peasants. Once the middle class becomes conservative, rural rebellion is still possible, but revolution is not.''
    \item Three levels of political participation:
    \begin{enumerate}
        \item Low: only the elites participate
        \item Medium: the aristocracy has some participation as well as portions of the upper-middle classes, but the ``masses'' remain generally excluded
        \item High: The entire population takes part in the political process
    \end{enumerate}
    \item \begin{center}
    \includegraphics[width=10cm]{"Reading Notes/Reading List/Huntington 1968/fig1.png"}\\
    Figure 1
    \end{center}
    \item ``As political participation increases, the complexity, autonomy, adaptability, and coherence of the society's political institutions must also increase if political stability is to be maintained.''
    \item \underline{Praetorian}: low institutionalization, high participation (this is mapped by Dahl as generally close to inclusive hegemony, which are rare). Lack of formal means of decision-making or adjudication -- 11social forces confront each other nakedly."  \underline{Civic}: high institutionalization, low participation; stable and recognizable patterns of institutional authority. High cost of authority limits resources that may be employed in politics and how groups may influence institutional processes
    \item Military intervention in politics is a reflection of underdeveloped institutions: the current system cannot appropriately satisfy the demand for participation; this is equivalent to how church, tribes, and other social groups become hyperinvolved in politics
    \item Revolution is much more rapid than reform and follows when there is not sufficient institutionalization to handle modernization -- instead, the increase in political participation spins out of control and requires almost a total overhaul of existing institutions. Conversely, reform occurs when the appropriate institutional capital is in place to manage these shifting needs, desires, and definitions of the selectorate; it is slower and more peaceful, but generally results in changes in institutions rather than total reconstruction
    \item Reform needs to be balanced: it must be fast enough to work and ``keep the dogs at bay'', but slow enough to remain under control. This requires strong, rather than weak, governments
\end{itemize}


\subsection{Important Specific Factual Claims}

\begin{itemize}[leftmargin=*, itemsep=4pt]
    \item Communist governments are highly effective at governing and establishing legitimate authority: ``Men may, of course, have order without liberty, but they cannot have liberty without order. Authority has to exist before it can be limited, and it is authority that is in scarce supply...''
    \item ``Modernization in practice always involves change and usually the disintegration of a traditional political system, but it does not necessarily involve significant movement toward a modern political system.'' At this time, devloping world had significant broad modernization, but not modernization of political systems -- just degradation, which led towards totalitarianism
    \item By drawing further conciousness to existing groups/social ties, modernization usually increases short-term conflicts and often results in fundamentalist religious movements. Creates group prejudice, which in turn creates group conflict
    \item Education and literacy are specifically tied to political stability... ``It is not the absence of modeernity but the efforts to achieve it which produce political disorder.''
    \begin{itemize}
        \item Counterpoint: ``In numerous underdeveloped countries the standard of living in the major cities is three or four times that prevalent in the countryside, yet the cities are often the centers of instability and violence while the rural areas remain quiet and stable. Political extremism is also typically stronger in the wealthier than in the poorer areas.''
    \end{itemize}
    \item quoting from Conroe: ``The higher the rate of change toward modernity, the greater the political instability, measured statically or dynamically.''
    \item ``At one extreme, some measure of economic growth is necessary to make instability possible. The simple poverty thesis falls down because people who are really poor are too poor for politics and too poor for protest.''
    \item ``Social mobilization is much more destabilizing than economic development.''
    \item Agrarian systems hinder mobility: limited opportunities for advancement within the agricultural sector
    \item ``The parties which at first are the leeches on the bureaucracy in the end become the bark protecting it from more destructive locusts of clique and family.''
    \item Modernization often increases gaps (including political) between urban and rural areas -- rural areas remain politically underdeveloped, but also with lesser need for enhanced political development (hence, why they generally remain peaceful). Also tends to create social disunity between these two cohorts that must be resolved to avert violence and unrest
    \item ``The sources of instability in a modernizing society are seldom in its poorest or most backward areas; they are almost always in the most advanced sectors of the economy''
    \item ``The extent to which the instability of the city manifests itself depends upon the effectiveness and the legitimacy of the political institutions of the society. Urban instability is minor but universal. Rural instability is major but avoidable.''
    \item ``Almost all praetorian oligarchies eventually evolve into radical praetorian systems.''
\end{itemize}

\subsection{Connections}

\begin{itemize}
    \item Many discussions of economic development and democracy. See Lipset 1959, A\&R, North, Iversen, etc.
    \item Ziblatt 2017: Huntington adds that there must also be a conservative movement within the middle class (Ziblatt focuses primarily on elites, but these are not mutually exclusive at all and in fact may be greatly supported by commonality)
    \item See Levy 2006 for discussion of tangible modernization processes / new goals on the state in advanced stages of development
\end{itemize}

\subsection{My Critiques}

\begin{itemize}
    \item Huntington is highly focused on institutional order rather than individual liberalism. This isn't necessarily a bad thing, but: while he warns of overly-burdensome institutions, he may underplay the potential costs of this institutionalization. Namely, he seems to suggest a rather monotonic relationship between state capacity and state quality that neither (A) satisfies contemporary normative claims nor (B) explain sustained instability in authoritarian regimes. Empirically, the work is heavily focused on the ``stability'' of Communism and the Soviet Union, which we now know to be untrue.
    \item Very specific procedural claims: X, then Y, then Z. Often there are at least some cases which support his transitions, but there are just as many which do not. For example, discussion of rapid urbanization, then overthrow of the rural elites by the urbanites, then partnership of urbanites and rural residents, is not an implausible pathway, but may be overly generalized here
    \item Not sure if his support of complexity is sufficiently nuanced. Do complex institutions really provide unilateral gains, or can decreased complexity better support high participation? Furthermore -- if high participation brings lower-educated cohorts into the political sphere, would highly complex systems become increasingly problematic? Could be a populism/demagoguery issue as well -- that which is not understood by the masses can be interpreted as animous
\end{itemize}


% ============================================================
\part{Iversen and Soskice 2019}
\begin{center}
Iversen, Torben, and David Soskice. \textit{Democracy and Prosperity: Reinventing Capitalism through a Turbulent Century}. Princeton: Princeton University Press, 2019.
\end{center}

\section{My Notes}

\subsection{Main Idea}



%--------------------------------------------------
\subsection{General Notes}
%--------------------------------------------------

\begin{itemize}[leftmargin=*, itemsep=4pt]
    \item ``As we see it, the advanced capitalist democracies, for all their instability and social problems not least at present, have been remarkably resilient and effective over this whole period. What we want to understand in the book is what mechanisms have driven that resilience and effectiveness over this long century.'' Creating a framework of capitalism within advanced economies; the core analyzed outcome is democratic resilience.
    \item Paints an opposing system to prominent claims that market capitalism has ``beat'' democracy, instead saying that democracy has evolved to encompass markets but been almost strengthened through it (VERY similar to the Levy 2006 argument)
    \item ``In our framework theory advanced capitalism is driven by the advanced democratic nation-state: democracy drives advanced capitalism. And in this process the autonomy of the advanced nation-state has increased even as globalization and mutual dependence have risen.'' Three elements of the reinforcing relationship between democracy and capitalism:
    \begin{enumerate}
        \item The role of government is central to facilitating effective capitalism and maximally free markets... ``Capitalism left to itself will hardly choose to operate in a competitive environment.''
        \item Parties to be electable have a reputation for managing advanced capitalism effectively
        \item Advanced capital is geographically embedded in the advanced nation-state rather than footloose (considering human and information capital rather than raw resources, for example)
    \end{enumerate}
    \item New cleavage between knowledge economy and low-productivity labor markets, yet claims it to not be a fundamental democratic threat -- largely because of the maintained prominence of the middle class and the consistent need for parties to secure its vote
    \item Democracy creates the conditions necessary for advanced markets to thrive
    \item Central issue regards the presence/nature of the median voter: are they within the middle class or the lower class? If in the middle class, then democracy and capitalism become self-sustaining and symbiotic, as the median voter is likely to support capitalistic endeavors. (Can also be fueled by the \emph{expectation} that one's children will have the opportunity to benefit through capitalistic endeavors)
    \begin{itemize}
        \item Further assumes that this is generally the case in advanced democracies: ``In our model a central component is that the large skilled workforces of the advanced sectors of the economy, and the aspirational voters who seek to join these workforces, have interests aligned with the promotion and success of advanced capitalist sectors, and are generally decisive voters.''
    \end{itemize}
\end{itemize}


\subsection{Important Specific Factual Claims}

\begin{itemize}[leftmargin=*, itemsep=4pt]
    \item End with a chapter about AI, though book was written in 2019 -- good discussion of emerging fronts
    \item Does NOT model capital as mobile, but rather as sticky: claims that this geographical rooting is an important 
    \item Populist sentiments arise out of geographic inequality: development of capital in one area does not ultimately lead to that capital being dispersed, creating ``losers'' in development
\end{itemize}


\subsection{Connections}

\begin{itemize}
    \item Against Boix 2003 that capital is mobile and that capital mobility provides institutional consolidation. Very clear connection here
    \item Significant footing in median voter theorm (also present in Muller \& Strom 1999)
\end{itemize}


\subsection{My Critiques}

\begin{itemize}
    \item
\end{itemize}


% ============================================================
\part{Kalyvas 2006}
\begin{center}
Kalyvas, Stathis N. \textit{The Logic of Violence in Civil War}. New York: Cambridge University Press, 2006.
\end{center}

\section{My Notes}

\subsection{Main Idea}



%--------------------------------------------------
\subsection{General Notes}
%--------------------------------------------------

\begin{itemize}[leftmargin=*, itemsep=4pt]
    \item Violence in civil war is a joint product of two series of interactions: those between actors at local versus central levels, and those between combatants and noncombatants
    \item Civil war presents an opportunity for local groups to pin their own policy or ideological goals onto the movement and in particular to use denunciation to impact otherwise extraneous concerns
    \item Violence can be indiscriminate or selective, but selective violence requires information that is costly/difficult to organize at the central unit and instead must be gathered through local networks (especially those of denunciation). As a result, civil war requires some form of alliance or cooperation between centralized actors who need information and local actors who need coercive forces. (This is the driving mechanism of what differentiates  selective versus indiscriminate violence -- the willingness of the state to endorse local yet extraneous conflicts and provide protection to denouncers encourages such denunciation, which allows for selective violence.)
    \item Violence is driven not by ethnicity / ideology but by the ability to achieve control over an area and the control already in place. Indiscriminate violence appears where control is low (defectors cannot be individually identified) while selective violence appears where control is moderate (defectors can be identified but are not dissuaded from defecting in the first place)
    \item Indiscriminate violence negates any incentive for individuals to join or support the perpetrators and is subsequently counterproductive: if you cannot guarantee the safety (relatively) of an actor who supports, then you cannot convince them to support if they otherwise would not
    \item Five zones of territorial control (reflected over a center point):
    \begin{itemize}
        \item 1 \& 5: total control by a single party
        \item 2 \& 4: partial (but dominant) control by one party
        \item 3: parity and contestation; expectation of maximal violence, particularly indiscriminate
    \end{itemize}
\end{itemize}


\subsection{Important Specific Factual Claims}

\begin{itemize}[leftmargin=*, itemsep=4pt]
    \item Unfolding identities and alliances through civil wars are often endogenous and not necessarily reflective of previous violence or cleavages
\end{itemize}


\subsection{Connections}

\begin{itemize}
    \item North \& Weingast as a commitment problem -- if you cannot create trust between the people and the state, then the state cannot sufficiently utilize the people as actors/agents of the state
\end{itemize}


\subsection{My Critiques}

\begin{itemize}
    \item
\end{itemize}


% ============================================================
\part{Levy 2006}
\begin{center}
Levy, Jonah D., ed. \textit{The State after Statism: New State Activities in the Age of Liberalization}. Cambridge, MA: Harvard University Press, 2006.
\end{center}

\section{My Notes}

\subsection{Main Idea}



%--------------------------------------------------
\subsection{General Notes}
%--------------------------------------------------

\begin{itemize}[leftmargin=*, itemsep=4pt]
    \item From preface: the role of the state has drastically changed in advanced democracies, but the presence of the state (and all it entails) is still very much there. How did these goals change, and what does the state look like in this modern construction?
    \item Two main conclusions from questioning new state missions across advanced democracies:
    \begin{enumerate}
        \item ``Contemporary developments in the economy, technology, society, and politics are not just challenging existing state activities but also creating demands and opportunities for new state intervention.''
        \item ``The spread of a liberal, market logic has been an especially powerful prod to state activism. In short, the move to the market has been, at the same time, a move to and through the state.''
    \end{enumerate}
    \item Shonfield: modern capitalism as ``the changing balance of public and private power,'' characterized by growth of public intervention in five areas (all quotes):
    \begin{enumerate}
        \item Management of the economic system
        \item Keynesian demand stimulus, the welfare state, and expanded public expenditures
        \item The ``taming'' of the ``violence of the market'' in the private sector, thanks to government regulation, planning, and inter-firm cooperation
        \item Policies to increase innovation and worker training
        \item The pursuit of intellectual coherence through long-term planning
    \end{enumerate}
    \item Broad shift in state activities from ``market-steering'' to ``market-supporting,'' with specific focuses (in this volume) on regulation, welfare, and labor policies
\end{itemize}


\subsection{Important Specific Factual Claims}

\begin{itemize}[leftmargin=*, itemsep=4pt]
    \item
\end{itemize}


\subsection{Connections}

\begin{itemize}[leftmargin=*, itemsep=4pt]
    \item Transition of institutional goals is almost directly linked to Huntington's theory of modernization -- emergence of new goals and institutions to achieve those goals
\end{itemize}


\subsection{My Critiques}

\begin{itemize}[leftmargin=*, itemsep=4pt]
    \item
\end{itemize}


% ============================================================
\part{Lijphart 1999}
\begin{center}
Lijphart, Arend. \textit{Patterns of Democracy: Government Forms and Performance in Thirty-Six Countries}. New Haven: Yale University Press, 1999.
\end{center}

\section{My Notes}

\subsection{Main Idea}



%--------------------------------------------------
\subsection{General Notes}
%--------------------------------------------------

\begin{itemize}[leftmargin=*, itemsep=4pt]
    \item Primary breakdown of majoritarian versus consensual government: government \emph{by} the people (representation in office-holding) versus government \emph{for} the people (reflection of representation in policymaking). Majoritarianism seeks to benefit the majority as explicitly as possible; consensus model seeks to include \emph{as many people as possible}. Exclusive/adversarial vs. inclusive/compromising; Kaiser 1997 terms it ``negotiation democracy''
    \item Two dimensions of democratic characterization:
    \begin{itemize}
        \item Executives-parties dimension
        \begin{enumerate}
            \item \underline{Concentration of executive power in single-party cabinets}
            \begin{itemize}
                \item Lowell's (1896) axiom posits that coalition cabinets are much weaker and less stable than single-party cabinets: ``The larger the number of discordant groups that form the majority, the harder the task of pleasing them all, and the more feeble and unstable the position of the cabinet.''
                \item Most appropriate metric of \# of parties seems to be the ``effective number of parties,'' calculated as 
                \begin{align*}
                    N &= \frac{1}{\sum s_i^2}
                \end{align*}
                where $s_i$ is the proportion of seats held by the $i$th party. Depends not only on number of parties, but on their shared relative strength: if heavily imbalanced, $N$ will decrease. Yet, this does not fully account for alliances / coalitions, whereby there may be some -- but not complete -- alignment between parties represented by different vote shares
                \item Dispersion of parties can be highlighted amongst seven issue dimensions:
                \begin{enumerate}
                    \item Socioeconomic
                    \item Religious
                    \item Cultural-ethnic
                    \item Urban-rural
                    \begin{itemize}
                        \item General shift in agrarian parties to enhance appeal towards suburban and rural voters as well, particularly as urbanization has decreased their voting share
                    \end{itemize}
                    \item Regime support
                    \item Foreign policy
                    \item Postmaterialist issues
                \end{enumerate} 
                However, not every democratic system is sensitive to each dimension, and in fact there is some range here: Lijphart characterizes Finland as embodying 3.5 dimensions (counting moderate engagement as one half), while Barbados has only 0.5 dimensions. Still, decently standardized range (not huge variance here). Potentially indicative of an ``exhaustion'' issue, whereby voters can truly only hold so many opinions? Or is this the result of party consolidation such that parties can't effective embody more than a handful of these dimensions without becoming far too isoalting?
                \item \emph{Positive relationship between \# of issue dimensions and effective number of parties: parties can only ``absorb'' so much of the policy space}. Taagepera and Grofman (1985) suggest that this can be estimated as $N = I + 1$ where $I$ is the number of issue dimensions
            \end{itemize}
            \item \underline{Dominance of executive over legislature}
            \begin{itemize}
                \item Various forms of cabinets, by majoritarianism usually supports a minimal winning coalition while consensus politics target oversized coalitions (which usually are not stable). General framework:
                \begin{center}
                    \includegraphics[width = 10cm]{"Reading Notes/Reading List/Lijphart 1999/fig 6.1.jpg"}
                \end{center}
                \item Many theories tend to point back to a MWC framework; Lijphart sees this as oversimplifying and ignoring evidence of oversized (and minority) coalitions
                \item Strong negative correlation between effective number of parties and the percentage of either MWC or single-party cabinets (this may be tautological though?)
                \item Previous literature often uses cabinet durability as a metric of executive dominance; Lijphart agrees that it can be considered in part but disagrees that it should be used as a primary metric alone
                \item Strong positive relationship between executive dominance and the percentage of cabinets which are either MWC or one-party
            \end{itemize}
            \item \underline{\# of parties (2 vs. many)}
            \item \underline{Majoritarian versus proportional representation}
            \begin{itemize}
                \item Implied district threshold: $T = \frac{75}{M + 1}$ where $M$ is the average district magnitude
            \end{itemize}
            \item \underline{Pluralist interests versus ``corporatist'' centralized interest group systems}
        \end{enumerate}
        \item Federal-unitary dimension
        \begin{enumerate}
            \item \underline{Unitary (centralized) versus federal (decentralized)}
            \item \underline{Unicameral versus bicameral legislature}
            \item \underline{Flexibility of constitutions (accessibility of amendments)}
            \item \underline{Presence/absence of judicial review}
            \item \underline{(In)dependency of central banks from the executive}
        \end{enumerate}
        Primary possible explanation of clustering within these dimensions: Goodin (1996)'s distinction of ``collective agency \& shared responsibility'' with divided agencies and responsibilities
    \end{itemize}
    \item Central focus on the ``Westminster'' Model, equated to the majoritarian framework (and particularly that model within the United Kingdom). Defined by dominance of a usually single-party cabinet ruling with a very slight majority over a primary opposition party, elected by a majoritarian system. Unitary governmental system with a unicameral legislature. Broad constitutional flexibility (lack of a centralized UK constitution) with no judicial review, and a politically-dependent central bank.
    \item Strongly supports consensus democracies: claims they outperform majoritarian democracies in most metrics and claim few costs. Particularly identifies consensus as being highly effective in plural societies (those with many or significant cleavages) by reducing electoral stakes and protecting minorities
\end{itemize}


\subsection{Important Specific Factual Claims}

\begin{itemize}[leftmargin=*, itemsep=4pt]
    \item Relative internal homogeneity can reduce the downside of majoritarianism: even though exclusion from the ruling regime may not prove representative to the minority, if their interests are quite similar to those in the majority, then presumably the policy outcomes should be relatively congruent regardless. That is, even though government \emph{by} the people doesn't hold, government \emph{for} the people does.
\end{itemize}


\subsection{Connections}

\begin{itemize}
    \item Main findings from preceding book \emph{Democracies} (1984):
    \begin{enumerate}
        \item Democratic institutions can be measured on a scale from majoritarianism to consensus
        \item These two institutional structures form two clusters
        \item A 2-D conceptual map of democracy can be used to place these democracies
    \end{enumerate}
    \item References Huntington 1991 in suggesting a ``two-turnover test'': that a government can be recognized as a stable and consolidated democracy once a leadership transition $A \to B \to A$ occurs (not only can the opposition gain power, but they can also cede it back). Also uses Huntington's 1991 characterization of the three waves of democratization to identify eligible states
    \item Lewis (1965) argument of the incompatability of majoritarianism with democracy: ``to exclude the losing groups from participation in decision-making clearly violates the primary meaning of democracy.''
    \item Directly uses Dahl's polyarchy as a means of defining democracy and therefore selecting qualifying regimes
    \item \textcolor{red}{Primary difference from Powell is that Lijphart explicitly supports consensus democracy over majoritarianism: claims evidence shows that economic performance is not sacrified, but various other metrics of institutional success are enhanced under a proportional system}
    \item Boix 1999 uses Effective Electoral Threshold, similar methodological framework for measuring the ability of any party to gain meaningful representation
    \item Tsebelis 1999 builds on this but explains parties as partial solutions to the veto-player problem (this hypothesis suggests that a wide range of veto players -- and incongruency between them -- should negatively impact legislative capacity)
\end{itemize}


\subsection{My Critiques}

\begin{itemize}
    \item Opposition to MWC seems a bit contrived -- even though there are individual examples of oversized majorities, these are usually the case precisely because the conditions which drive MWC outcomes have not \emph{yet} come into effect. Namely, there is a high chance for collapse or consolidation later; that is, Lijphart is identifying outcomes, but not equilibrium outcomes.'
    \item Falls pretty bluntly within the comparative instinct to over-categorize... Created eight possible arrangements of legislative/executive tradeoffs and only three had any states (one of which only had one)
    \item Direction of causation in many of these relationships is unclear: does an increase in number of dimensions cause an increase in the number of parties, or vice versa? It is pretty clear that these can't be exclusively determined by the information given, but Lijphart seems to make some directional implications that are not necessarily substantiated in any meaningful sense
\end{itemize}


% ============================================================
\part{Lipset and Rokkan 1967}
\begin{center}
Lipset, Seymour Martin, and Stein Rokkan. ``Cleavage Structures, Party Systems, and Voter Alignments: An Introduction.'' In \textit{Party Systems and Voter Alignments: Cross-National Perspectives}, eds. Seymour Martin Lipset and Stein Rokkan. New York: Free Press, 1967.
\end{center}

\section{My Notes}

\subsection{Main Idea}



%--------------------------------------------------
\subsection{General Notes}
%--------------------------------------------------

\begin{itemize}[leftmargin=*, itemsep=4pt]
    \item Contemporary party divides are primarily the result of the consequences from two major historical transitions: the national revolution, and the industrial revolution. Party populations have been largely frozen around their primary-served groups following these transitions, demonstrating a great deal of stability with respect to their primary constitutent bases. This is largely determined by the order and combination of the above junctures and their effects, rather than by contemporary considerations. (\textcolor{red}{freezing hypothesis})
    \item Four cleavages generated by these revolutions, upon which party ``identities'' have largely been based:
    \begin{enumerate}
        \item Center vs. periphery
        \item State vs. church
        \item Land vs. industry
        \item Owner vs. worker
    \end{enumerate}
    \item Cleavages require three components in order to become politically consequential or relevant to political identity:
    \begin{enumerate}
        \item Empirical basis: there must be some sort of identifiable social division, a way to distinguish the cleavage itself and include/exclude individuals based on their characteristics
        \item Normative element: there must be some set of shared identity, ideals, or interests, such that those values can legitimate the division between the ``out''-groups
        \item Organizational expression: groups, parties, institutions, etc. which mobilize and represent the division (this is the difference between cleavages and plain social difference)
    \end{enumerate}
    \item A political movement must cross four thresholds to warrant full institutionalization as a party:
    \begin{enumerate}
        \item Legitimation: members can petition, criticize, and organize as an opposition group within the broader political Systems
        \item Incorporation: legislative suffrage is extended to the movement's supporters
        \item Representation: the movement can, in fact, win seats
        \item Majority power: the movement can plausibly compete for a position in the government, and is not a perpetual minority voice
    \end{enumerate}
    \item Cleavages can either be cross-cutting or reinforcing:
    \begin{itemize}
        \item Cross-cutting cleavages occur when some class or group defined by a cleavage is itself divided by another cleavage, resulting in greater party fragmentation and cross-pressurization of voters
        \item Reinforcing cleavages occur where cleavages are largely aligned, such that subgroups are not further split and group identity is reinforced. This tends to create polarization along a smaller number of much deeper lines
    \end{itemize}
\end{itemize}


\subsection{Important Specific Factual Claims}

\begin{itemize}[leftmargin=*, itemsep=4pt]
    \item Use of an A/G/L/I spatial model, but somewhat ill-defined and difficult to work through...
    \item Parties have an incentive/mechanism for their own survival: once established, they will evolve and adapt such that some version of the institution itself can survive. This becomes a foundation of path dependency; it is difficult to change directions once the party has been established
\end{itemize}


\subsection{Connections}

\begin{itemize}
    \item General connection on group identity -- see Anderson 1983
    \item Many of the legitimation mechanisms seem to tie back to Powell 2000 and a question of majoritarian governance -- a proportional approach may allow for a group to become politically relevant with much greater accessibility
    \item Another connection to Dahl: is the political system inclusive, and is the opposition meaningfully part of the political process
    \item Similar arguments to Moore 1966's characterization of political systems as a reflection of historical orders of wealth and welfare, though along generally different approaches
    \item See Diermeier and Vlaicu 2011 for additional discussion of how like-minded groups might be able to collaborate on effective policymaking even without formal party "whips"
    \item See Posner (2004) and Liesner \& Fearon (1997) -- ethnic identities/cleavages not necessarily the drivingfactors themselves in conflict or tension, but rather the underlying means through which salient divisions/tension can arise. Needs to be some contributing factor (usually the impact of such cleavages on developing inequality between groups)
\end{itemize}


\subsection{My Critiques}

\begin{itemize}
    \item Is this still true? Factually, it seems like it may have fallen apart quite a bit, particularly as left-of-center parties have aligned around more educated and wealthier voters. It is also a tad difficult to truly consider existing party conditions in 1920s Europe given the presence (and fate) of the Weimar Republic
\end{itemize}


% ============================================================
\part{Mill 1843}
\begin{center}
Mill, John Stuart. \textit{A System of Logic, Ratiocinative and Inductive}. Book III, Chapter VIII: Of the Four Methods of Experimental Inquiry
London: John W. Parker, 1843.
\end{center}

\section{My Notes}

\subsection{Main Idea}



%--------------------------------------------------
\subsection{General Notes}
%--------------------------------------------------

\begin{itemize}[leftmargin=*, itemsep=4pt]
    \item Five primary approaches to inquiry:
    \begin{itemize}
        \item \underline{Method of Agreement}: If a phenomenon occurs multiple times, and in every case of its occurence there exists only one common factor, then that factor is the cause. (does not guarantee unique causation)
        \item \underline{Method of Difference}: If a phenomenon occurs in one scenario and does not occur in the other, and the two scenarios are identical other than one factor, then that factor is causal. (satisfies necessary \& sufficient, \textcolor{red}{is strongest method available})
        \item \underline{Joint Method of Agreement and Difference}: If multiple instances of a phenomenon have one common factor, while multiple instances where it does not occur are only tied by the absence of that common factor, then that factor is causal
        \item \underline{Method of Residues}: If $A + B + C \rightarrow E_1 + E_2 + E_3$ is observed, and $A \rightarrow E_1$ and $B \rightarrow E_2$ is known, then $C \rightarrow E_3$. (requires prior knowledge of some of the causal chains)
        \item \underline{Method of Concomitant Variations}: If variation in one phenomenon is tied to variation in another phenomenon,then there exists a causal chain between them. (conflates correlation and causation, quite problematic)
    \end{itemize}
    \item Broad focus on eliminating non-causal factors (this also ties to positivity)
\end{itemize}


\subsection{Important Specific Factual Claims}

\begin{itemize}[leftmargin=*, itemsep=4pt]
    \item
\end{itemize}


\subsection{Connections}

\begin{itemize}
    \item See Lipset 1959: ``A deviant case, considered within a context which marshals the evidence on all relevant cases, often may actually strengthen the basic hypothesis if an intensive study of it reveals the special conditions which prevented the usual relationship from appearing.'' pg. 70
\end{itemize}


\subsection{My Critiques}

\begin{itemize}
    \item
\end{itemize}


% ============================================================
\part{Moore 1966}
\begin{center}
Moore, Barrington, Jr. \textit{Social Origins of Dictatorship and Democracy: Lord and Peasant in the Making of the Modern World}. Boston: Beacon Press, 1966.
\end{center}

\section{My Notes}

\subsection{Main Idea}



%--------------------------------------------------
\subsection{General Notes}
%--------------------------------------------------

\begin{itemize}[leftmargin=*, itemsep=4pt]
    \item Central question: ``We seek to understand the role of the landed upper classes and the peasants in the bourgeois revolutions leading to capitalist democracy, the abortive bourgeois revolutions leading to fascism, and the peasant revolutions leading to communism.''
    \item English liberalization largely intertwined with changing (increasing) value of land. In late 1500s-early 1600s, land values rapidly increased, and lords began to stake claim to additional lands (legally or extrajudiciously) that had previously been communally used by the peasantry. Transition from collective/customary ruling to private ownership, largely driven by the \underline{yeomen}: smaller-scale capitalists who began to commercialize, but were not as wealthy as the ``true'' aristocratic class. The yeomen critically were not direct collaborators with the Crown -- and in fact, began to push back as their wealth grew. The origins of modern Parliament subsequently began as a ``committee of landlords'' to satisfy / enfranchise this group post-Civil War.
    \item Development of commercial agriculture (largely in line with this above change in the value of land and a movement towards exploitation rather than sustenance) is a driving factor in the development of a liberal society: it is largely the primary instigator of private markets which drives broad social change and allows for / needs new institutions (Huntington?)
    \item One core factor in development: whether or not the aristocracy was dependent on the crown/state or on the market. If dependent on the market, it was much harder to fully upend, and they would have a much better chance of securing political enfranchisement. If dependent on the state, much more ripe for more violent responses
    \item Difficult to have peaceful transition to enfranchisement when there is a large and effectively powerful peasant class. If the peasant class can be neutralized in the years leading up to a significant change in agricultural output, then the odds of peace are much higher -- but, if the peasantry is a significant force, their drastically increased output can lead to a need for violent revolution. This is largely in line with Huntington: if economic wellbeing expands rapidly but political institutions have not meaningfully developed, there will be unrest and violence until an equilibrium is reached
    \item \textbf{``No bourgeois, no democracy.''}
    \item Four pre-conditions for democracy:
    \begin{enumerate}
        \item Development of a balance of power between the crown/central authority and the landed aristocracy that prevents excessive royal absolutism;
        \item A turn toward commercial agriculture that weakens rather than reinforces repressive labor controls;
        \item Weakening or destruction of the landed aristocracy as an independent political class (or its fusion with the bourgeoisie);
        \item prevention of an aristocratic-peasant coalition against the bourgeoisie/commercial classes.
    \end{enumerate}
    \item Boix 2003 agrees w/ Moore: ``This model of political transition... matches and gives internal consistency to Moore's insight about the antidemocratic nature of landlords; at the same time, it accounts for Moore's intuition about the bourgeoisie's benign role in the process of extending universal suffrage."
\end{itemize}


\subsection{Important Specific Factual Claims}

\begin{itemize}[leftmargin=*, itemsep=4pt]
    \item On English political development in light of the expanded wool trade: ``Under the pressure of circumstances, the medieval notion of judging economic actions according to their contribution to the health of the social organism began to collapse. Men ceased to see the agrarian problem as a question of finding the beset method of supporting people on the land and began to perceive it as the best way of investing capital in the land''
    \item British Civil War strengthened the land-owning class while the French Revolution weakened it: in Britain, this transfer of power to the wealthy laid the framework for future (peaceful) transitions but required the violence of the civil war to set it up. Yet, this empowered the large farming tenants to make jumps in agrarian production, and they ultimately drove this economic transition (capital was not provided by the landlords)
    \item Three possible pathways to modernization, contingent on the above commercialization of agriculture:
    \begin{itemize}
        \item Bourgeois-democratic route (England, France, USA)
        \item Capitalist-reactionary route to fascism (Germany, Japan)
        \item Communist route (Russia, China)
    \end{itemize}
\end{itemize}


\subsection{My Critiques}

\begin{itemize}
    \item
\end{itemize}


% ============================================================
\part{Muller and Strom 1999}
\begin{center}
M\"uller, Wolfgang C., and Kaare Str\o m, eds. \textit{Policy, Office, or Votes? How Political Parties in Western Europe Make Hard Decisions}. Cambridge: Cambridge University Press, 1999.
\end{center}

\section{My Notes}

\subsection{Main Idea}



%--------------------------------------------------
\subsection{General Notes}
%--------------------------------------------------

\begin{itemize}[leftmargin=*, itemsep=4pt]
    \item Parties have three possible goals: shaping policy, holding office, or gaining votes. These goals are not necessarily complimentary (and indeed are often in conflict) -- cannot be maximized simultaneously
    \item Consider a party as three leading actors: its elected officials, its central leadership, and its members/activists, all of whome have different priorities regarding the above goals
\end{itemize}


\subsection{Important Specific Factual Claims}

\begin{itemize}[leftmargin=*, itemsep=4pt]
    \item Suggests that refusing office is a means of pursuing policy or vote-maximization goals
\end{itemize}


\subsection{Connections}

\begin{itemize}
    \item Harmel \& Janda 1994 -- discusses party goals as varying between policy, office, or votes. Also discuss a fourth motivation: democracy-seeking
\end{itemize}


\subsection{My Critiques}

\begin{itemize}
    \item
\end{itemize}


% ============================================================
\part{Nalepa 2010}
\begin{center}
Nalepa. \textit{Skeletons in the Closet: Transitional Justice in Post-Communist Europe}. New York: Cambridge University Press, 2010.
\end{center}

\section{My Notes}

\subsection{Main Idea}



%--------------------------------------------------
\subsection{General Notes}
%--------------------------------------------------

\begin{itemize}[leftmargin=*, itemsep=4pt]
    \item Differentiating why some post-USSR democracies pursued lustration, the discovery/vetting/exposure of those who were collaborators with the secret polise, while others did not (or only did to varying extents)
    \item Game theoretic conclusion -- lustration depends on politicians' private incentives to oppose lustration on the grounds of their own previous transgressions (i.e., the skeletons in the closet)
    \item Those without skeletons stand to gain normatively and politically from pursuign lustration, including in future political races or discussions
    \item Previous literature suggests this should just be in line with a party's previous support for communism; work shows that this is not the case
\end{itemize}


\subsection{Important Specific Factual Claims}

\begin{itemize}[leftmargin=*, itemsep=4pt]
    \item Motivating instance is the bargaining over transition of power, primarily as it stands to significantly influence whether or not documents/materials are kept or destroyed
\end{itemize}


\subsection{Connections}

\begin{itemize}[leftmargin=*, itemsep=4pt]
    \item
\end{itemize}


\subsection{My Critiques}

\begin{itemize}[leftmargin=*, itemsep=4pt]
    \item
\end{itemize}


% ============================================================
\part{North 1990}
\begin{center}
North, Douglass C. \textit{Institutions, Institutional Change and Economic Performance}. Cambridge: Cambridge University Press, 1990.
\end{center}

\section{My Notes}

\subsection{Main Idea}



%--------------------------------------------------
\subsection{General Notes}
%--------------------------------------------------

\begin{itemize}[leftmargin=*, itemsep=4pt]
    \item Institutions are the rules of the game: ``humanly devised constraints that structure political, economic, and social interactions.'' Organizations are the players of the game
    \item Institutions exist to reduce transaction costs and facilitate a more efficient market solution
    \item Biggest non-obvious contribution is the commitment issue: the state must gain power while committing to those who give it power that it will not use that power for abuse or extraction. Economic growth happens precisely when institutions (a) minimize transaction costs and (b) maximize property rights and constrain the sovereign
    \item Institutions are path-dependent, as they facilitate increasing returns (i.e., are more effective as time goes on) -- status quo reinforcement
    \item Most critical metric is an institutional environment's ability to adapt to new settings and experiment with various adaptations
    \item Institutions are sometimes changed by those who stand to personally improve with respect to that change -- institutional opportunism
\end{itemize}


\subsection{Important Specific Factual Claims}

\begin{itemize}[leftmargin=*, itemsep=4pt]
    \item Turn towards institutionalism as a framework of theory and methodology
\end{itemize}


\subsection{Connections}

\begin{itemize}
    \item
\end{itemize}


\subsection{My Critiques}

\begin{itemize}
    \item
\end{itemize}


% ============================================================
\part{Olson 1965}
\begin{center}
Olson, Mancur. \textit{The Logic of Collective Action: Public Goods and the Theory of Groups}. Boston: Harvard University Press, 1965.
\end{center}

\section{My Notes}

\subsection{Main Idea}



%--------------------------------------------------
\subsection{General Notes}
%--------------------------------------------------

\begin{itemize}[leftmargin=*, itemsep=4pt]
    \item Collective action problem exists -- rational actors will not support a collective good if they have an incentive to deviate / shirk
    \item Group size as a key variable (the larger the group, the bigger the free-riding problem)
    \item Effective groups often instead provide some degree of private services bundled with common resources, such that there is some distinct individual incentive to contribute. Collective ends should be by-products of operations
    \item Public goods are non-excludable (cannot prohibit access for non-contributors) and non-rival (one person's consumption does not reduce availability)
    \item A group is ``privileged'' if it is rational for at least one member to bear the cost alone (still can introduce a coordination problem though)
    \item Must provide some form of selective incentives where receipt (positive or negative) can be targeted based on contribution
\end{itemize}


\subsection{Important Specific Factual Claims}

\begin{itemize}[leftmargin=*, itemsep=4pt]
    \item ``Exploitation of the great by the small'' -- those with the largest stakes will usually endure the largest costs (disproportionately) as they are most incentivized to maintain provision
\end{itemize}


\subsection{Connections}

\begin{itemize}[leftmargin=*, itemsep=4pt]
    \item
\end{itemize}


\subsection{My Critiques}

\begin{itemize}[leftmargin=*, itemsep=4pt]
    \item
\end{itemize}


% ============================================================
\part{Ostrom 1990}
\begin{center}
Ostrom, Elinor. \textit{Governing the Commons: The Evolution of Institutions for Collective Action}. London: Cambridge University Press, 1990.
\end{center}

\section{My Notes}

\subsection{Main Idea}



%--------------------------------------------------
\subsection{General Notes}
%--------------------------------------------------

\begin{itemize}[leftmargin=*, itemsep=4pt]
    \item Tragedy of the commons occurs without supporting institutions; when CPRs are properly managed by means of institutional development, the tragedy can be avoided. Leviathan solution is not optimal for this use case
    \item Best solution to the commons is the establishment of self-organized and self-governed institutions -- those who are subject to the resource and its usage should create them to manage
    \item Communication, monitoring, customization, and enforcement are all necessary traits of collective action solutions
    \item Institutional construction is a costly endeavor, creating an additional potential collective action problem -- easier to let someone else take the initiative of institutional design
    \item Potentially actually most appropriate as a lens of institutional endogeneity rather than as a substantive application?
\end{itemize}


\subsection{Important Specific Factual Claims}

\begin{itemize}[leftmargin=*, itemsep=4pt]
    \item
\end{itemize}


\subsection{Connections}

\begin{itemize}
    \item
\end{itemize}


\subsection{My Critiques}

\begin{itemize}
    \item
\end{itemize}


% ============================================================
\part{Powell 2000}
\begin{center}
Powell, G. Bingham Jr. \textit{Elections as Instruments of Democracy: Majoritarian and Proportional Visions}. New Haven: Yale University Press, 2000.
\end{center}

\section{My Notes}

\subsection{Main Idea}

There are two primary frameworks (and institutional structures) through which elections can define democratic governance: majoritarian, and proportionalistic. In a majoritarian approach, the government is bound to directly reflect the interest of the party receiving a majority of seats (though this may not always imply a majority of votes), and provides a rather ``all-or-nothing'' approach to representation. Conversely, in a proportionalistic system, the government ideally reflects the proportionate makeup of policy preferences and perspectives throughout the voting population, and delivers policy influence in line with those proportions. Oftentimes, the proportional approach provides better voter-representative congruency and fairer policy representation as compared to the majoritarian lens. It further tends to more effectively protect minority interests, and in particular can help balance policymaking. Yet, majoritarian governance implicitly improves voter-policy congruence. That is, while proportionalism can deliver a legislature that most effectively reflects the people, it may not empower such legislatures to be meaningfully effective in any broadly-improving way.



%--------------------------------------------------
\subsection{General Notes}
%--------------------------------------------------

\begin{itemize}[leftmargin=*, itemsep=4pt]
    \item Primary distinction between democratic approaches: \emph{control} over policymakers or \emph{influence} over policy. Each corresponds to a certain ``grand vision'' of electoral accountability:
    \begin{itemize}
        \item \underline{Majoritarian}: The majority should have a near-absolute right to rule, and policy should proceed in accordance with the majority's wishes. Focuses on a relationship of power between voters and their elected officials, whereby voters become ``principals'' who seek to either reward or control their representatives. Consistent with the ``control'' narrative. Centralization of power follows from this, and the primary role of the citizen is to enforce consequences during elections.
        \item \underline{Proportionalism}: The will of each group should be reflected in policy in accordance with the relative size of that group. Focuses on a relationship of influence whereby policy is swayed in a degree proportionate to the relative power of that group, subsequently protecting minority welfare at the expense of direct control over the political system. Consistent with the ``influence'' narrative. Considers elections to be clumsy/noisy signals, and prioritizes universal input without giving an oversized benefit to the majority group
    \end{itemize}
    \item Voters' roles can be roughly perceived as either retrospective or prospective. If retrospective, more aligned with majoritarianism: the voter either punishes or rewards incumbents with their vote (or lack thereof), based on the observed quality of their tenure. If prospective, more aligned with proportionalism: the voter seeks credible commitments from candidates and has some expectation of their performance.
    \begin{center}
        \includegraphics[width=10cm]{"Reading Notes/Reading List/Powell 2000/powell fig 1.1.jpeg"}
    \end{center}
    \item The result of this broader differentiation is two metrics of electoral success: (1) responsiveness to elections in choosing policymakers, and (2) representational congruence between the policymakers' positions and the citizens' preferences. Majoritarianism \emph{should} do better with (1), and proportionalism better with (2).
    \item Election rules in SMDP systems generally limit the number of seats and subsequent equilibria features limited numbers of viable candidates based off of those seats (see Duverger). This can create a significant disproportion in representation amongst majoritarian systems: a 49\% share may ultimately be entirely unrepresented in a legislature. High effective thresholds -- explicit or systemic vote share requirements in order for a party to gain representation -- benefit a majoritarian perspective. As the effective threshold increases, the expected vote-seat disproportionality increases as well.
    \item ``If the decision rules encourage or require the government parties to negotiate with independent groups in the assembly, especially those representing voters who supported different parties, I want to characterize them as closer to the proportional vision. If the decision rules encourage the government parties to control the executive and the assembly without check or division, I characterize them as majoritarian.''
    \item ``Influence of the opposition'' from Strom -- largely based on the role and importance of committees. Presence is important, of course, but further considerations: do the committees have policymaking power? Is there any veto opportunity for minority parties within committees? Are the committees highly specialized such that mainstream ideological backgrounds of certain parties can be effectively cast aside? Committees as a primary metric of decentralization
    \item Though constitutional design and legislative rules can change over time, their change in democracies is usually not particularly drastic or quick, for a number of reasons including institutional safeguards. But, an additional -- and potentially strongest -- consideration: an incumbent usually has won office through some means of benefitting from the system in place, and therefore has minimal incentive to change it. This is a more self-serving equilibrium outcome than a Dahl ``trust in the opposition'' framework, and rather suggests tautologically that the system usually has worked for whoever is currently in power.
    \item A core component of restrospective voting requires that the citizens are able to perceive who should be held accountable for what decisions. Though this may be obvious in some more direct settings -- and majoritarian systems provide a great deal of transparency -- it becomes a bigger issue with coalition governments. Powell finds that the extent to which voters treat an in-government party as responsible for policy throughout their incumbency is contingent on their position within the coalition as well as whether or not they ruled individually or with other parties.
    \item Party mandates stem from a majoritarian perspective whereby the government's winning of a majority of votes/seats/etc. provides it with both the ability and the responsibility to fully deliver on its visions, even occasionally with disregard for institutional safeguards or compromises. Mandates ``assume a single decisive choice'' and generally simplify policy into one dimension; this is potentially an over-simplification that blurs their ability to govern appropriately on less-salient concerns.
    \item Control through mandates requires two stages:
    \begin{enumerate}
        \item ``The voter needs to be able to identify the prospective future governors and have some idea of what they will do if elected. This condition is necessary if the voters are to make good prospective choices.'' Calls them ``\emph{identifiable prospective governments}"
        \item ``The outcome of the election should bring into office a coherent government committed to policies that correspond to the voters' anticipations and capable of carrying them out. Calls them ``\emph{a responsively formed working majority after the election}''
    \end{enumerate}
    \item ``In the nineteenth-century context having a voice in the legislature may have seemed sufficient. But the emergence of extremely cohesive legislative parties and the dominance of many parliaments by their executives imply to those concerned about them the desirability of giving minorities some greater weight in policy making than merely the opportunity to be heard in (largely irrelevant) legislative debate.'' That is, PR protects minority interests to a better extent than majoritarianism. Yet, simple resolutions to this often come at a great normative cost: empowering minorities to enact vetos is likely inherently undemocratic, and any pathway leaves plenty of opportunities to be abused.
    \begin{itemize}
        \item ``Given that citizens' preferences are appropriately articulated by the parties they choose to support in the election, there remain two conditions: that each party win fair representation in the assembly and that it have fair influence in policy making."
    \end{itemize}
    \item Finds general support for PR particularly wrt effective authorized representation: tends to provide unilaterally-improved representation compared to MAJ systems. The work of course considers this conclusion to be nuanced, but an important consideration in constitutional design.
    \item Responsiveness is not the only metric of successful democracy and in fact may often be in conflict with other measures of success. Consider abuse of majoritarianism, lack of delegation to public policy experts, etc. Two hypotheses for responsiveness:
    \begin{enumerate}
        \item ``The majoritarian vision relies on elections featuring few competitors and identifiable future governments during the election and majority governments that thoroughly control policy making after the election.... A single party wins a majority of votes, controls the legislature, and makes policy. Governments can claim voter mandates for their election promises and are clearly responsible for their actions.''
        \item ``The proportional vision implies its own hypotheses about responsiveness. This vision stresses the superiority of all citizens and the need for rules that fairly reflect the choices of citizens in the composition of the legislature. The legislature suffers from neither overweighting of some parties at the expense of others nor forced cohabitation within individual parties or preelection coalitions that can freeze bargaining opportunities.''
        
        \begin{center}
        \includegraphics[width=10cm]{"Reading Notes/Reading List/Powell 2000/Powell Fig 6.1.jpg"}
        \end{center}
    \end{enumerate}
    \item Party vote and policy-making share are generally much more aligned under a PR system than majoritarian (see below section for more details on distortion). 
    \item The decisive stage of citizen preferences in the majoritarian model is the election itself -- yet, in the proportionalistic view, it is multistage. Citizens are represented once in the individual election, and then again in coalition bargaining between parties if no outright majority is achieved.
    \item A proportionality approach \emph{can} improve voter-policy congruence; a majoritarian approach \emph{does} improve it. ``In situations in which a single party or coalition of parties forms a government and must maintain tight party discipline, which empirically is the case in almost all parliamentary systems, the govern,ment might be expected to make policies that correspond not to the legislative median but to its own internal distribution.''
    \item ``Why do majoritarian processes so often fail? One answer is that despite the apparenet robustness of the theoretical connections, correspondence depends in fact on achieving some very specific conditions and, importantly, that it is the nature of the majoritarian processes that fairly small deviations from these lead to large, undesirable consequences. A second answer is that it frequently is asking too much of elections that in a complex world they bear the full, decisive burden of connecting voters and policymakers as the majoritarian vision presumes.''
\end{itemize}


\subsection{Important Specific Factual Claims}

\begin{itemize}[leftmargin=*, itemsep=4pt]
    \item Making retrospective judgments requires some clarity about specifically \emph{who} made \emph{what} judgments. Requires both some degree of transparency and a sense of consistency from amongst the citizens, who in turn need to accurately perceive the dynamic at hand. This adds an additional dynamic to be considered: \underline{clarity of responsibility}
    \item Focuses on two alternative methods of connection between citizen preferences and the formation of policy coalitions (and subsequently policy): differentiating factor is whether or not a voter's vote is sufficient to fully express their preferences. This may not be the case precisely because voters are presented with an exclusive menu, rather than empowered to fully express their specific preferences over a continuous domain.
    \item Possible veto institutions to restrict a national assembly's power:
    \begin{enumerate}
        \item Independent executives (main example is USA)
        \item Separate legislative chamber with independenet selection base and veto powers
        \item Federal systems: regional checks on centralized power
        \item Judicial review (think Vanberg) -- court system (usually independent) has some oversight over legislative process
    \end{enumerate}
    \item Majoritarianism provides a degree of quantitative stability in election result shifts: a tiny change in vote margin is unlikely to result in a regime change in a majoritarian system, while a proportionalist system may be much more volatile to such changes
    \item Party mandates ``seems to have appeared at the end of the nineteenth century as the first mass political parties, usually Socialist parties, were mobilizing new electorates in Western Europe.... They argued, moreover, that, having made policy commitments to the electorate, they were both authorized and obligated to carry out those policies once in office. This formulation translated the older concept of majoritarian or populist democracy into terms that took explicit advantage of the enfranchisement of mass electorates and the emergency of mass political parties."
    \item Majoritarian systems are more likely to allow for severe distortion whereby a majority of votes does not result in a majority of seats (by clumping votes within districts, etc.) because of its all-or-nothing nature. However, this sort of distortion -- though possible -- is much less common within PR systems, as there are fewer opportunities for broad deviations. Unlikely to see minority-vote-holders become majority-seat-holders in a proportional system (that is, outright -- does not consider inclusion in a coalition government, which is its own consideration).
    \item Two types of responsiveness failures:
    \begin{enumerate}
        \item ``Those in which failing to meet one ideal draws the election outcome closer to the alternative ideal''
        \item ``Unmitigated failures, in which distance from one ideal does not draw the outcome toward the alternative''
    \end{enumerate}
\end{itemize}


\subsection{Connections}

\begin{itemize}
    \item Figure 1.2 (pg. 15) should link directly to Dahl: begin with citizen preferences/beliefs, flow to policy
    \item Many references to (and in agreement with) Lijphart 1984, who differentiates two approaches to democratic government: either the government should be responsive to the majority (Westminster approach), or should be responsive to ``as many people as possible."
    \item ``The effective representation of voters for opposition parties is always less than that of voters for governing parties, but the degree of difference varies, depending on the bargaining resources of those parties.'' $\rightarrow$ very reminiscent of Dahl's polyarchy argument: is the opposition competitive, and is the political system competitive as a whole?
    \item Connection to Carey (2009) and Larreguy \& Raffler (2025) regarding different models of accountability, etc. Likely a tie between majoritarianism and the rise of social media being against it, but the pros of majoritarianism are therefore degraded
    \item Boix 1999: mainstream parties use the transition from majoritarianism to PR as a means of protectionism when a new party (often socialist) emerges -- but, this requires the ability to consolidate around a given establishment party, which can be difficult / impossible if the ruling parties are relatively comparable in size and influence
    \item Diermeier \& Vlaicu 2011: discuss majoritarianism as a reflection of pivotal politics, and further expand on how it can be sustained even in the absence of formal party organizations
    \item Gandhi \& Przeworski 2007: further expansion of the "mandated right to rule" within authoritarian systems, whereby autocrats allow for some dissent (particularly by means of inclusive legislatures) such that their power can actually be further consolidated due to the additional credible threats they have lent
\end{itemize}


\subsection{My Critiques}

\begin{itemize}
    \item Doesn't seem to particularly explore the geographic aspect here: proportionality makes it much harder to have individual district representatives who are accountable for certain aspects, etc, without making the number of delegates huge (or the number of districts small). It seems possible that a PR approach may simply not be effective in smaller political bodies where granularity cannot be achieved while holding direct accountability.
    \item Discusses the inherent stability of majorities, given that few fall apart post-election -- but, isn't this somewhat tautological? Aren't these systems unlikely to collapse simply because there is virtually no institutional means by which this can happen? Perhaps this also follows from a Duverger's application, but in a two-party legislature, it is presumably rare (though not impossible) that a schism in one party results in absorption by the other...
\end{itemize}


% ============================================================
\part{Przeworski et al 2000}
\begin{center}
Przeworski, Adam, Michael E. Alvarez, Jose Antonio Cheibub, and Fernando Limongi. \textit{Democracy and Development: Political Institutions and Well-Being in the World, 1950--1990}. New York: Cambridge University Press, 2000.
\end{center}

\section{My Notes}

\subsection{Main Idea}



%--------------------------------------------------
\subsection{General Notes}
%---------------------------------------------anne-----

\begin{itemize}[leftmargin=*, itemsep=4pt]
    \item ``The world around us is not generated randomly. The world we observe is a result of actions of people pursuing their ends. It nurtures successes and eliminates failures. These processes are systematic, and so are their effects. The search for randomness is futile; we must proceed differently. The question is how.'' \\
    ``To cope with the non-random nature of our observations, we revert to counterfactuals. To evaluate institutions, to assess their impact, we must ask what the outcomes would have been if other institutions had been observed under the same conditions. And if their selection is not random, these institutions are not observable under the same conditions. Hence, a reoucrse to counterfactuals is unavoidable.''
    \item Contestation requires three virtues (largely defined by Przeworski's \emph{a system in which parties lose elections}):
    \begin{enumerate}
        \item Ex-ante uncertainty: some positive probability that some member of the ruling government will lose an election
        \item Ex-post irreversibility: whoever wins the election will assuredly assume office
        \item Repeatability: whoever gains office will not rewrite the electoral institutions such that they guarantee a future win
    \end{enumerate}
    \item Though there are certainly examples (especially in 20th century) of autocracies experiencing significant economic growth, generally not controversial that democracy seems to be correlated with development -- the question is the direction of causality; it is often claimed (here argued incorrectly so) that democracy itself causes this economic growth.
    \item Distinction between two core causal mechanisms: emergence and survival. Data shows that while democracies are not necessarily more likely to emerge as economic development increases, when economic hardship is introduced, a democratic system is much more likely to fail than an authoritarian system. Subsequently, there emerges a survivorship bias, as those democracies which experience poor economic results are effectively no longer categorized as democracies. Conversely, dictatorship survival is less tied to economic development -- ``dictatorships die under the entire range of economic conditions.'' pg. 111
    \item ``The most important lesson we have learned is that wealthy countries ttend to be democratic not because democracies emerge as a consequence of economic development under dictatorships but because, however they emerge, democracies are much more likely to survive in affluent societies. We find it difficult to explain why dictatorships die and democracies emerge. Although we are willing to believe that economic development may open the possibility for transition to democracy, even when the conditions for democracy are ripe, the outcomes of political conflicts are indeterminate. Hence, we failed to detect any thresholds of development that would make the emergence of democracy predictable. In sum, modernization theory appears to have little, if any, explanatory power.'' pg. 137
\end{itemize}


\subsection{Important Specific Factual Claims}

\begin{itemize}[leftmargin=*, itemsep=4pt]
    \item ``A mypoic public has to be protected from itself; economic reforms must be carried out before the people have time to react. This is still a popular belief.'' pg. 4
    \item Good section on methodology in the introduction -- discusses how measuring welfare via income is inherently incomplete, as well as the difficulty of democracies turning into dictatorships when the economy is bad (makes it very hard to continuously code data and will cause an apparent shock to the democratic results)
\end{itemize}


\subsection{Connections}

\begin{itemize}[leftmargin=*, itemsep=4pt]
    \item Strong tie between the above discussion of counterfactuals and Mill's system of inductive logic \& reasoning: application of Mill's approach to comparative institutions (may need to verify that they are exactly aligned though)
    \item Heavy tie to Dahl's polyarchy: Przeworski et al focus on contestation as the primary distinction between democracies and non-democracies, whereby democracies feature legitimately contested political power and influence (most notably the executive and the legislature)
    \item Very aligned with Huntington 1968 in that a lack of political development can be destabilizing in the face of economic growth: if political institutions are not sufficiently developed to endure greater participation and interest (namely, democracies should be very well-equipped to do so), then we should expect to see instability; here, instability is tied to survival of the regime as a whole
\end{itemize}


\subsection{My Critiques}

\begin{itemize}[leftmargin=*, itemsep=4pt]
    \item
\end{itemize}


% ============================================================
\part{Putnam 1993}
\begin{center}
Putnam, Robert D. \textit{Making Democracy Work: Civic Traditions in Modern Italy}. Princeton: Princeton University Press, 1993.
\end{center}

\section{My Notes}

\subsection{Main Idea}



%--------------------------------------------------
\subsection{General Notes}
%--------------------------------------------------

\begin{itemize}[leftmargin=*, itemsep=4pt]
    \item Central question: ``What are the conditions for creating strong, responsive, effective representative institutions?'' Comparison of regional Italian bureaucracies as an extended case, exploiting the simultaneous creation of fifteen new regional governments in the 1970s: homogenous constitutional authority, but drastically different policy results. Argument is that differing social and political contexts lead to these results
    \item Transformation of existing metrics of institutional performance: we do not merely need the government to make decisions, but we further need them to \emph{act} on those decsions. Agreement itself (a social choice theoretic outcome) is not sufficient to demonstrate institutional success. Model is:
    
    \begin{align*}
        \text{Societal demands} & \implies \text{Political interaction} \\
        & \implies \text{Government} \\
        & \implies \text{Policy choice} \\
        & \implies \text{Implementation}
    \end{align*}
    ``A high-performance democratic institution must be both responsive and effective: sensitive to the demands of its constituents and effective in using limited resources to address those demands.''
    \item Core transition in institutional authority came with the ability to establish their own patronage system: here called ``safeguarding public morals,'' the new institutions (regional) were ultimately allowed to create their own licensing, permitting, and regulatory efforts that could exploit (and be exploited by) local actors.
    \item With regionalization came a new generation of regional elites, defined by a ``passing of the torch'' to those with lesser name recognition and hereditary status and instead self-made wealth and status
    \item Any metric of government performance should satisfy the following:
    \begin{enumerate}
        \item \underline{Comprehensive}: must consider the wide range of potential roles of government, and act accordingly
        \item \underline{Internally consistent}: must consider the fact that governments can be measured along different frameworks and dimensions, and must approximately hold that rank
        \item \underline{Reliable}: should be generalizable rather than highly volatile: can/should change as performance does, but should not be super-overly-reactive in the absence of meaningful reasons to redirect
        \item \underline{Correspond to the objectives and evaluations of the institution's protagonists and constituents}: must actually reflect the needs/desires/metrics of the people subjected to them
    \end{enumerate}
    \item One driving theory of differential efficacy stems from the civic context: the ``civic virtue'' of the people, as historical normative theory would call it. This was often in contrast with republicanism, but had re-swept Western thought in the mid-20th century. 11As the proportion of nonvirtuous citizens increases significantly, the ability of liberal societies to function successfully will progressively diminish."
    \item Some possible dimensions of civic capacity:
    \begin{itemize}
        \item Civic engagement: are citizens active in the political process? Not necessarily a 1:1 here, as action can by non-virtuous or otherwise fueled by solely self-interest. Does not defeat a collective action problem -- most citizens still look for short-term maximization of the immediate family
        \item Political equality: do citizens interact as equals, with themselves and with their government officials?
        \item Solidatiry, trust, and tolerance: what do the bonds between citizens look like? Are they strong even across heterogenous groups, identities, and backgrounds?
        \item Associations: are there distinctive social structures and norms? Are there civil and political organizations, particularly those founded by citizens rather than imposed by the government? (This is a bit of an Ostromian thought as well, regarding endogenous institutions) Also civic institutions that are entirely apolitical -- think later Putnam ``Bowling Alone''
    \end{itemize}
    \item A significant factor is social trust: citizens in civic communities trust each other and expect fair dealings, and similarly trust the government's ability to achieve high standards in exchange for compliance (a heightened social contract). ``Collective life in the civic regions is eased by the expectation that others will probably follow the rules. Knowing that others will, \emph{you} are more likely to go along, too, thus fulfilling \emph{their} expectations.''
    \item In the absence of such systems of trust and reciprocal enforcement of collective action problems, less civic communities turn towards forces of coercion (i.e., the police). This is the only alternative solution to the Leviathan. ``In the absence of solidarity and self-discipline, hierarchy and force provide the only alternative to anarchy.'' This also results in a paradox whereby more individualistic societies are actually more state-dependent for enforcing order.
    \item Heightened social trust allowed for financial growth and significant expansion of trade; in turn, the prosperity these returned built even further social trust (and trust in the civic institutions). A sign/portion of this was the rise of credit and loanmaking.
    \item Positive relationship between civic traditions and civic community; directionality is a bit hard to establish, but plausible that having an established set of traditions allows for social coherence which in turn lays the ground for civic community.
    \item \emph{``Economics does not predict civics, but civics does predict economics, better indeed than economics itself.''}
    \item Social capital -- and ``rotating credit associations'' -- can be a solution to the collective action problem, by means of building trust in one another and subsequently allowing for higher-risk activities and investments. Rotating credit association members buy into a fund which is distributed later to one member. Subsequently, social capital ``is the key to making democracy work.''
\end{itemize}


\subsection{Important Specific Factual Claims}

\begin{itemize}[leftmargin=*, itemsep=4pt]
    \item Arturo Israel: It is easier to build a road than to build an organization to maintain that road
    \item Pre-regional reformation in the 1970s, public policy in Italy was almost entirely the result of centralized oversight and public affairs, mediated by heavy patronage influences. This centralization was maintained throughout various iterations of the national Italian Government
    \item Distinction between North and South in Italy: the North was more concerned by central threats to authority, while the South was more concerned by central threats to funding. North favored ``horizontal" collective action across a wide front while South favored ``vertical'' action through private petitions, etc.
    \item ``The first test for any new political institution is that it must engage the aspirations and harness the ambitions of serious politicos.''
    \item Broad depolarization of Italian politics (at least at the local level) through and beyond this period of reformation; shift of local politics from ideological extremes to pragmatic policymaking. Much of this may be explained by institutional socialization, as the new institutional structure resulted in more moderates holding power and facilitated moderating interactions between ideologically opposed actors
    \item Ties back to early medieval times, whereby self-governance (in some form) arose out of serfdom in the north more than in the south of Italy -- ultimately became the communes, first as voluntary organizations and later as actual governance. Still had a sense of self-enforcement, where illicit activity was met with ostracization. Beginning of the oaths of mutual security; rise of the notary, which became a significant institution in organizing civic life
    \item ``Industrial districts'' with high rates of production/manufacturing etc. Characterized by a mix of vigorous competition between firms to innovate, but cooperation to establish common administrative/legal/financing/research outcomes. This creates a mix of institutional adaptiveness and economic strength that ultimately fueled a great deal of economic development in late-20th-century Italy.
    \item Social capital as a ``moral resource'': must be used in order to secure its presence and power, and actually decreases in strength if not exercised. Thus, perpetuates greater activity (when appropriately set)
\end{itemize}


\subsection{Connections}

\begin{itemize}
    \item Intense social and economic development in post-WW2 Italy sets a similar context a la Huntington 1968, particularly as political institutions did \textit{not} keep up with this level of development. 
    \item Directly references Huntington's concepts of questioned authority / development of autonomy away from previously-tribalized structures: regional autonomy in Italy as a prime example of that, as people began to question patronage/heritage-based systems, etc.
    \item Another reference to ``input'' vs. ``output'' of government a la Dahl 1972: government input is the processes and institutions which exist and are constructed to govern, and output are the actual laws and policies that come with them
    \item Other scholars (Edmund Burke, Bariel Almond, back to Cicero) cite political consensus as a foundation for good government; Putnam disagrees. Rather, civic communities / networks allow for dissent to be appropriately handled throughout society without driving signficant schisms. Evidence in Italy does not suggest that institutional success was correlated with political homogeneity
    
\end{itemize}


\subsection{My Critiques}

\begin{itemize}
    \item None. I love Putnam and he is near and dear to my heart
\end{itemize}


% ============================================================
\part{Scott 1985}
\begin{center}
Scott, James C. \textit{Weapons of the Weak: Everyday Forms of Peasant Resistance}. Bethany, Connecticut: Brevis Press, 1985.
\end{center}

\section{My Notes}

\subsection{Main Idea}



%--------------------------------------------------
\subsection{General Notes}
%--------------------------------------------------

\begin{itemize}[leftmargin=*, itemsep=4pt]
    \item While literature has previously focused on collective peasant actions en masse, argues that the driving factor of political development is decentralized, lower-level everyday resistance (also particularly targeting indirect resistance)
    \item Everyday resistance can include intentional inefficiency, destruction of property, false compliance, feigned ignorance, foot-dragging... all construct a degree of protection against exploitative bodies
    \item Everyday resistance is rationally viable particularly because the costs are low and there exists no meaningful collective action barrier to overcome
    \item Resistance is often symbolic in some part -- actions are heavily judged by meaning and ideological intentions
    \item As landowners and wealthier farmers gained representation/influence/wealth, common peasants (evidence here is from Sedaka in Malaysia) responded with everyday rather than organized resistance in an effort to preserve previous economic structure, rights, situation, etc.
    \item Characterizes much peasant activity as at a level below organized revolt, but instead efforts to slowly squeeze against higher classes or other influential bodies/groups
\end{itemize}


\subsection{Important Specific Factual Claims}

\begin{itemize}[leftmargin=*, itemsep=4pt]
    \item Uses Green Revolution as a natural experiment
\end{itemize}


\subsection{Connections}

\begin{itemize}
    \item
\end{itemize}


\subsection{My Critiques}

\begin{itemize}
    \item
\end{itemize}


% ============================================================
\part{Selznick 1949}
\begin{center}
Selznick, Philip. \textit{TVA and the Grass Roots: A Study in the Sociology of Formal Organization}. Berkeley: University of California Press, 1949.
\end{center}

\section{My Notes}

\subsection{Main Idea}



%--------------------------------------------------
\subsection{General Notes}
%--------------------------------------------------

\begin{itemize}[leftmargin=*, itemsep=4pt]
    \item Organizations are faced with an existential gap between what they wish to do and what they are able to accomplish; subsequently, the natural outome is ``a reconciliation between the desire and the ability.''
    \item Organizations ``take on the aspects of a social organism,'' with the following consequences:
    \begin{enumerate}
        \item Decentralization without harming organizational goals is possible only when a common set of beliefs are deeply held / indoctrinated amongst multiple levels of leadership \emph{below the very top} -- that is, organizational goals and values are internalized such that a greater degree of discretion is possible
        \item Consistent doctrine allows for technologically-specialized experts to bind together when they may otherwise be quite drawn to their individual groups / professional interests. 
        \item Institutions gain a tendency to defend its own functions and existence. ``These defensive activities are aided when a set of beliefs is so fashioned as at once to fortify the special needs or interests of the organization and to provide an aura of disinterestedness under which formal discussion may be pursued.''
    \end{enumerate}
    \item Institutions must adapt to their environments, but: at least in the case of the TVA, adaptation to local institutions was much more significant (both in presence and necessity) than adaptation to local populations. Though local populations may very well hold strong beliefs, it is the institutions which represent those populations who may actually impact the route to implementation, positively or negatively. (For the TVA, this apporoach represented a pragmatic rather than normative effort -- grassroots organizations allowed for effective progress, not simply a reflection of ``good governance.'')
    \item Eight inherent dilemmas regarding administrative/policymaking discretion (all quotes):
    \begin{enumerate}
        \item There is a dilemma of doctrine and commitment, or of the abstract and concrete.
        \item There is also the dilemma of consent and conformance, or of selective decision and total involvement.
        \item So far as discretion is delegated, bifurcation of policy and administration is reinforced.
        \item Theories about government become preempted by social forces.
        \item Emphasis on existing institutions as democratic instruments may wed the agency to the status quo.
        \item Decision at the grass roots may be inhibited by the system of national pressure groups.
        \item Commitment to existing agencies may shape and inhibit policy in unanticipated ways.
        \item Existing agencies inhibit a direct approach to the local citizenry.
    \end{enumerate} 
\end{itemize}


\subsection{Important Specific Factual Claims}

\begin{itemize}[leftmargin=*, itemsep=4pt]
    \item ``Organizations established in a stable context, with assured futures, do not feel the same urgencies as those born of turmoil and set down to fend for themselves in undefined ways among institutions fearful and resistant. As in individuals, insecurity summons ideological reinforcements.'' Broader points of institutional threat -- argument possibly that it could result in better institutions under such stressors, though not entirely clear
\end{itemize}


\subsection{Connections}

\begin{itemize}
    \item Discussions of cooptation and decentralization are largely aligned with Huber \& Shipan -- ideological congruency allows for greater discretion, which in turn is productive for facilitating grassroots organization. A broader tension though arises when there is not ideological congruency and therefore cannot be sufficient discretion given to lower levels of authority in order to promote a grassroots approach.
    \item Institutions adapting to existing institutions rather than existing populations -- very much a Lijphart/Powell viewpoint of representative democracy (majoritarian) whereby elected officials are granted a mandate
\end{itemize}


\subsection{My Critiques}

\begin{itemize}
    \item
\end{itemize}


% ============================================================
\part{Stokes 2013}
\begin{center}
Stokes, Susan C., Thad Dunning, Marcelo Nazareno, and Valeria Brusco. \textit{Brokers, Voters, and Clientelism: The Puzzle of Distributive Politics}. New York: Cambridge University Press, 2013.
\end{center}

\section{My Notes}

\subsection{Main Idea}



%--------------------------------------------------
\subsection{General Notes}
%--------------------------------------------------

\begin{itemize}[leftmargin=*, itemsep=4pt]
    \item Often an incentive to target poor voters, but a long history of electoral exchanges:
    \begin{quote}
        ``If there's a fire on Ninth, Tenth, or Eleventh Avenue, for example, any hour of the day or night, I'm usually there with some of my election district captains as soon as the fire-engines. If a family is burned out... I just get quarters for them, buy clothes for them if their clothes were burned up, and fix them up till they get things runnin' again. It's philanthropy, but it's politics, too -- mighty good politics. Who can tell how many votes one of these fires bring me? The poor are the most grateful people in the world."
    \end{quote}
    \item Core breakdown of distributive politics: distinctions between following vs. not following public, binding rules; attempts to influence voters vs. attempts to hold voters to account
    \begin{center}
        \includegraphics[width=14cm]{"Reading Notes/Reading List/Stokes 2013/fig1.jpg"}
    \end{center}
    \item A significant distinction (and in fact, the defining characteristic between programmatic and non-programmatic politics) is whether or not the rules (if followed) are public: is there some backdoor exchange, or is it a clear mechanism which simply funnels certain distribution? Note that this requires (or at least calls for) an institutional consideration: if there are not formal institutions shaping the redistributive process, then presumably there are also not actual rules publicly accessible and shaping said distribution. Per SCOTUS 1982:
    \begin{quote}
        ``We have never insisted that the franchise be exercised without tain of individual benefit; indeed, our tradition of political pluralism is partly predicated on the expectation that voters will pursue their individual good through the political process, and that the summation of these individual pursuits will further the collective welfare. So long as the hoped-for personal benefit is to be achieved through the normal processes of government, and not through some private arrangement, it always has been, and remains, a reputable basis upon which to cast one's ballot."
    \end{quote}
    \item Quid-pro-quo exchanges under clientilism are a normative threat: ``These exchanges seem to violate the free action or autonomy of voters. Even if we accept that voters are never fully autonomous and always come under the influence of some other actor -- parents, co-workers, ``opinion makers,'' or party leaders -- still the image of the voter being held to account for his or her choice is disquieting.''
    \item Clientelism requires an extraordinary amount of information regarding individual voters: their needs, their wants, their political leanings, and their ultimate political actions. The solution to this information problem is ``brokers,'' large quantitites of individuals employed by political machines to both gather and operate on this information. The broker is also tasked with providing specific benefits and actually implementing the ``get" side of distributive politica.Yet, by introducing brokers as contracted with political parties or centralized machines, we create a principal-agent tradeoff whereby there are potentially incentives to shirk or exploit. In short: vote buying works because of, and through, the political broker.
    \item Brokers have a clear incentive to target swing voters, and parties are not able to identify exactly who might be a swing voter versus safely support/opposition. Clientelism allows for these swing voters to be targeted by means of the broker's information and local expertise. Yet, empirically, this is not always true: in fact, swing voters are often under-represented in clientelism, while safe supporters tend to reap benefits at a disproportionate rate.
\end{itemize}


\subsection{Important Specific Factual Claims}

\begin{itemize}[leftmargin=*, itemsep=4pt]
    \item
\end{itemize}


\subsection{Connections}

\begin{itemize}
    \item See Nichter 2008: extension of Stokes 2005 data on Argentina, but suggests a further model for clientilism that is much less dependent on accurate monitoring of individuals. Core extension: what if we cannot assume that brokers have a reasonably accurate perception of individual votes, and instead can only observe turnout? Provides an alternative answer to the puzzle that Stokes et al present, as the model ultimately suggests that clientilism should indeed target ideologically-aligned individuals. Also directly discusses turnout buying, one of the foundational types identified by Stokes. Stokes is a bit skeptical about this, so paints a divide
    \item Müller 1997: Discusses a transition away from party operatives/oganizations to instead surveys and mass media, particularly regarding socialist parties and mass movements in the mid-20th century. Presents a somewhat conflicting argument here (or at least, a different perspective in some cases) regarding whether or not the formal institution of the party itself is particularly relevant.
    \item Carey \& Shugart 1995: when should candidates campaign on party versus personal reputation? Presumably, clientelism occurs through drastically different means depending on this level -- whether the party as a whole can support it (what Stokes calls clientelism more broadly) or whether by the individual (turnout buying?)
    \item Crisp et al 2004: discussion of ``local public goods'' and other various forms of provision, particularly focusing on when legislators prioritize local or national concerns within their own actions
    \item Schmitter 1974 -- corporatism as a particular example of this, although the key difference is that with corporatism the benefitting body has some legitimate power and influence over policy decisions
\end{itemize}


\subsection{My Critiques}

\begin{itemize}
    \item Less of a critique and more of a side point -- the work doesn't fully consider pressuring voters to abstain, either via intimidation or simply through general persuasive means. It may not be possible to convince dissidents to vote in line with the party, but it may be more feasible to convince them not to vote at all... (Addressed but only very briefly)
\end{itemize}


% ============================================================
\part{Svolik 2012}
\begin{center}
Svolik, Milan W. \textit{The Politics of Authoritarian Rule}. Cambridge: Cambridge University Press, 2012.
\end{center}

\section{My Notes}

\subsection{Main Idea}



%--------------------------------------------------
\subsection{General Notes}
%--------------------------------------------------

\begin{itemize}[leftmargin=*, itemsep=4pt]
    \item Two central conflicts of authoritarian politics:
    \begin{enumerate}
        \item \underline{The problem of authoritarian control}: Balancing against the majority excluded from power; maintaining authority and legitimacy over the people. Think external threats
        \item \underline{The problem of authoritarian power-sharing}: Countering challenges from other authorities/institutions who have some degree of power or influence over the current political process. Think internal threats
    \end{enumerate}
    Considers most authoritarian institutions (including legislatures, etc.) to be attempts to solve either of these issues, or both.
    \item Because dictatorships lack independent oversight institutions, agreements between autocratic leaders cannot be enforced by law or other ``hard-institutional'' means -- rather, they must be maintained through some degree of interpersonal agreement and a subsequent equilibrium outcome based on soft institutions. \textbf{self enforcing} (A further extension of this is the role of violence in authoritarian regimes, which replaces much of the institutional process seen in democratic systems.)
    \item Four dimensions of dictatorship:
    \begin{enumerate}
        \item Military involement in politics
        \item Restrictions on political parties
        \item Concentration of power with the dictator
        \item Personalization of political interactions
    \end{enumerate}
    \item There exists inherent tension between the dictator and their allies, as the dictator could use his position to expand his authority and opportunity for rent-seeking; allies fear they could be eventually excluded from the selectorate. The balance to this is the credible threat of removal of the dictator by their allies. This yields two distinct forms of authoritarian power-sharing:
    \begin{enumerate}
        \item \underline{Contested autocracy}: The dictator's allies can pose a revolutionary threat to the dictator that is sufficiently credible in order to force concessions. (Likely characteristic of the USSR post-Stalin)
        \item \underline{Established autocrats}: The dictator's allies do not pose a credible revolutionary threat and therefore cannot force concessions. (Common in African dictatorships)
    \end{enumerate}
    \item Though difficult, it is possible to establish institutions which alleviate the power-sharing commitment problem. Such institutions must monitor compliance and allow for penalties upon noncompliance. Two types of primary institutions along this route:
    \begin{enumerate}
        \item Deliberative and decision making institutions: when discussions and deliberations occur with some degree of transparency, it is easier to maintain appropriate monitoring of the dictator's actions. Such institutions facilitate regular interaction between the dictator and his allies and accordingly yield this transparency (at least to some extent). This largely explains the presence of legitimate legislatures and legislative/executive councils in dictatorships -- they are not necessarily elected nor directly empowered, but their existence supports a power-sharing balance.
        \item Formal institutions are less ambiguous and more broadly known by those involved, allowing for the detection of any noncompliance with the power-sharing agreement
    \end{enumerate}
    \item Allies face a collective action problem in responding to noncompliance or other opportunistic affairs: if only some subset seeks to punish the ruler for noncompliance, they may easily be bested and subsequently severely punished. Rather, enforcement of noncompliance requires a broad range of involvement that may be difficult to achieve even in the face of clear reneging.
    \item Complementary relationship between repression and cooptation, rather than a distinct trade-off between the two: the goal is to coopt those individuals and institutions which may be either sympathetic to the dictator or sufficiently influential enough to warrant the cost of this coercion, while repressing those individuals and institutions who would not be susceptible to such ``buy-in'' (or would require an unjustifiable amount of effort to secure their participation)
\end{itemize}


\subsection{Important Specific Factual Claims}

\begin{itemize}[leftmargin=*, itemsep=4pt]
    \item Aristotle \emph{Politics} book 3: ``We have next to consider how many forms of government there are, and what they are.... Tyranny is a kind of monarchy which has in view the interest of the monarch only; oligarchy has in view the interest of the wealthy; democracy of the needy. None of them the common good of all.''
    \item Significant discussion of the definition of dictatorship; notes that it must first be evaluated as binary and then evaluated as a degree. Fundamentally implies that the legislature is not comprised of freely and fairly elected officials (and to a smaller extent, that the executive is not elected freely and fairly OR by the free \& fair legislature).
    \item A transition from contested to established autocracy is virtually equivalent to a transition from oligarchy to autocracy
\end{itemize}

\subsection*{Connections}

\begin{itemize}
    \item Fourth dimension is personalization of political interactions -- ties directly back to Stokes et al 2013 on clientelism; potentially ties back to Dahl's Polyarchy as a closed hegemony
    \item Definitely well-connected to Gandhi \& Przeworski (2007) in a general sense of ``why'' authoritarian institutions matter
\end{itemize}

\subsection{My Critiques}

\begin{itemize}
    \item
\end{itemize}


% ============================================================
\part{Weber 1930}
\begin{center}
Weber, Max. \textit{The Protestant Ethic and the Spirit of Capitalism}. Trans. Talcott Parsons. London: George Allen \& Unwin, 1930 [1904--1905].
\end{center}

\section{My Notes}

\subsection{Main Idea}



%--------------------------------------------------
\subsection{General Notes}
%--------------------------------------------------

\begin{itemize}[leftmargin=*, itemsep=4pt]
    \item Some similarity between Marx and Weber, in that both were seeking to explain modern economic development and consider the conditions it both required and created -- but, Weber took specific pieces of Marxist ideology while explicitly rejecting its total implementation/acquisition. Weber was particularly critical of some of the Marxist followers of this time and especially their revolutionary sentiments; instead sought a much more peaceful transition and to explain the organic development of the modern economy.
    \item Though desires for wealth acquisition are not novel, the contemporary West has become associated with ``the rational organisation of formally free labour.'' -- its ``routinised, calculated administration within continuously functioning enterprises.'' This implies both a disciplined labour force and the regularised investment of capital. (Quote from interpreter but not from Weber)
    \item At this time, economic development and wealth accumulation seemed to be highly correlated with the presence of Protestant communities: such societies were achieving much greater economic development than their non-Protestant counterparts; in particular, Catholics.
    \item Weber claims that the underlying connection between Protestantism and a developed work ethic -- particularly through the ``virtuous'' reinvestment of capital and pursuit of meaningful labor -- were particularly key components in the development of modern economies (and in the heterogeneity of such development). Everyday men sought to pursue meaningful work, which in turn led to the development of massive amounts of capital and development
    \item Protestantism encouraged individuals to pursue ``good works'' for the sake of labor rather than pleasure or joy; as a result, production and investment became intrinsically virtuous activities (think acts of creation)
    \item Key argument: this occurred as an unintended consequence; Protestantism discouraged the agglomeration of great amounts of wealth, but conversely created an environment ripe for productivity and investment of capital
\end{itemize}


\subsection{Important Specific Factual Claims}

\begin{itemize}[leftmargin=*, itemsep=4pt]
    \item In a context of predestination post-Calvin, people sought to prove externally that they were to be saved, and performed worldly ``good works'' as a demonstration of surety. Subsequently, fulfilling one's calling was no longer a means by which salvation could be achieved, but a signal to others that the individual was to be saved.
\end{itemize}

\subsection{Connections}

\begin{itemize}
    \item Lipset 1959: ``For the purpose of this paper, democracy (in a complex society) is defined as a political system which supplies regular constitutional opportunities for changing the governing officials. It is a social mechanism for the resolution of the problem of societal decision-making among conflicting interest groups which permits the largest possible part of the population to influence these decisions through their ability to choose among alternative contenders for political office. In large measure abstracted from the work of Joseph Schumpeter and Max Weber, this definition implies a number of specific conditions: (a) a ``political formula,'' a system of beliefs, legitimizing the democratic system and specifying the institutions -- parties, a free press, and so forth -- which are legitimized, i.e. accepted as proper by all; (b) one set of political leaders in office; and (c) one or more sets of leaders, out of office, who act as a legitimate opposition attempting to gain office.'' Talks quite a bit more about Weber as well
\end{itemize}

\subsection{My Critiques}

\begin{itemize}
    \item
\end{itemize}


% ============================================================
\part{Weber 1968}
\begin{center}
Weber, Max. ``The Nature of Charismatic Authority and Its Routinization'' and ``The Pure Types of Legitimate Authority.'' In \textit{Max Weber on Charisma and Institution Building: Selected Papers}, ed. S.\ N.\ Eisenstadt, pp.\ 46--65. Chicago: University of Chicago Press, 1968. Excerpted from \textit{The Theory of Social and Economic Organization}.
\end{center}

\section{My Notes}

\subsection{Main Idea}



%--------------------------------------------------
\subsection{General Notes}
%--------------------------------------------------

\begin{itemize}[leftmargin=*, itemsep=4pt]
    \item Stable domination or equilibrium order is \emph{Herrschaft}, which requires a claim to legitimacy. Three possible such claims:
    \begin{enumerate}
        \item Traditional authority -- norms, cultural rules, dynasties, etc.
        \item Charismatic authority -- devotion to a specific leader
        \item Legal-rational authority -- legitimacy vested in institutions or rules
    \end{enumerate}
    \item Charismatic authority is extra-institutional and is largely unstable; it cannot survive for an extended period of time and instead must either adapt or die. This adaptation is routinization -- charismatic authority must instead be shifted into embedded social and political structures, transforming charistmatic authority into either traditional or legal-rational authority
\end{itemize}


\subsection{Important Specific Factual Claims}

\begin{itemize}[leftmargin=*, itemsep=4pt]
    \item
\end{itemize}


\subsection{Connections}

\begin{itemize}[leftmargin=*, itemsep=4pt]
    \item Miller 1990: importance of charismatic leadership given that incentive structures cannot be perfectly constructed
\end{itemize}


\subsection{My Critiques}

\begin{itemize}[leftmargin=*, itemsep=4pt]
    \item
\end{itemize}


% ============================================================
\part{Weber 2013}
\begin{center}
Weber, Max. 2013. \textit{From Max Weber: Essays in Sociology} (eds.\ and trans.\ H.H. Gerth and C. Wright Mills). Hoboken: Routledge. pp.\ 77--83, excerpt from ``Politics as a Vocation.''
\end{center}

\section{My Notes}

\subsection{Main Idea}



%--------------------------------------------------
\subsection{General Notes}
%--------------------------------------------------

\begin{itemize}[leftmargin=*, itemsep=4pt]
    \item State as the monopoly of violence
    \item Tripartite theory of political authority suggests three origins of political authority
    \begin{enumerate}
        \item Charismatic
        \item Traditional
        \item Rational-legal
    \end{enumerate}
    \item Politics: ``Striving to share power or striving to influence the distribution of power, either among states or among groups within a state''
    \item Achieving the Herrschaft requires not only coercive capacity but also legitimacy
\end{itemize}


\subsection{Important Specific Factual Claims}

\begin{itemize}[leftmargin=*, itemsep=4pt]
    \item
\end{itemize}


\subsection{Connections}

\begin{itemize}[leftmargin=*, itemsep=4pt]
    \item
\end{itemize}


\subsection{My Critiques}

\begin{itemize}[leftmargin=*, itemsep=4pt]
    \item
\end{itemize}


% ============================================================
\part{Weiner 1962}
\begin{center}
Weiner, Myron. \textit{The Politics of Scarcity: Public Pressure and Political Response in India}. Chicago: University of Chicago Press, 1962.
\end{center}

\section{My Notes}

\subsection{Main Idea}



%--------------------------------------------------
\subsection{General Notes}
%--------------------------------------------------

\begin{itemize}[leftmargin=*, itemsep=4pt]
    \item ``At the heart of this study is a simple problem. In a free society, there is an increase in the freedom of choice and in the capacity to affect one's environment -- a freedom and a capacity exercised by people who voluntarily associate to influence the course of public policy. Freedom of association, freedom of assembly, a free press, and free speech are the freedoms maintained in a democratic society so that citizens may affect their government. On the other hand, voluntary associations would defeat their own purposes if they were to effectively paralyze government, destroy orderly rule, or threaten to Balkanize the political system."
    \item ``The nationalists were creating a Leviathan of aspirations'' -- might be wise to use this phrase a bit more broadly when considering the impact of rapid expansion and development on political outcomes and social cohesion
\end{itemize}


\subsection{Important Specific Factual Claims}

\begin{itemize}[leftmargin=*, itemsep=4pt]
    \item In early/mid-20th century: ``Indian leadrs saw their task as that of reducing scarcity, which was more acutely felt at the time of Independence than it had been fifty years before, not necessarily because there was greater poverty in India, but because popular aspirations had expanded much more rapidly than had the society's ability to fulfill public demands.'' Attempted to do so from a top-down approach: reducing scarcity should be the responsibility of the government
    \item As political groups began to disrupt the new social order and more mainstream attempts to satiate them were unsuccessful, Indian officials largely depended on criminal prosecutions and legal restraint/coercion, including those which specifically restrict aforementioned freedoms of speech and press. (State could arrest individuals in anticipation of a future public statement)
\end{itemize}


\subsection{Connections}

\begin{itemize}
    \item Seems to be a bit related to the Huntington argument, except instead of the underlying change being economic development, it is instead the development of freedom: political institutions must be sufficiently strong enough to handle new waves of freedom while still maintaining effective policies and scarcity management.
\end{itemize}


\subsection{My Critiques}

\begin{itemize}
    \item
\end{itemize}


% ============================================================
\part{Ziblatt 2017}
\begin{center}
Ziblatt, Daniel. \textit{Conservative Parties and the Birth of Democracy}. New York: Cambridge University Press, 2017.
\end{center}

\section{My Notes}

\subsection{Main Idea}



%--------------------------------------------------
\subsection{General Notes}
%--------------------------------------------------

\begin{itemize}[leftmargin=*, itemsep=4pt]
    \item Question: Why did historical elites and wealthy give away so much power and security to allow for democracy, rather than entirely preventing its arisal?
    \item Answer: ``The historical stabilization of European democracies required a long-rn shift, spanning the mid-1800s to the mid-1900s, in the type of political parties representing conservative social forces.... Conservatives discovered the mobilizing power of modern political parties, which allowed them to protect their own interests, in turn ironically smoothing the way for a more stable path of democratization. Political systems endured when conservatives organized their own political parties and embraced the virtues of pragmatic political action -- what nineteenth-century critics decried as political wirepulling.''
    \item Motivating argument: well-organized and effective conservative parties (representing the interests of the elites) was the alternative to revolution, and ultimately allowed greater security for the elites even as some non-violent loss was expected. Namely, if elites successfully organized these conservative movements, then they could use democracy to their own advantage, and enjoy tempered spoils while avoiding total eradication
    \item Methodologically notes the importance of examining long-run trends, namely by demonstrating the ability of cross-sectionally ``strong'' observed democracies to be unsuccessful over the long-run. Rather, evaluating historical trends allows us to more appropriately categorize different outcomes, rather than just trying to fit pattern-seeking to immediate cases
    \item Another mechanism to consider -- Conservative parties hindered overreach by political bodies other than the elites (sometimes the Church, but often the party or the state) and provided a mediating factor to expansion of the franchise (cooled things down, in a Huntington 1968 framework). But, when conservative parties are not equipped to do so -- notably, when met with extremist or populist movements that forego these traditional conservative political values -- this balance becomes upset, and can often swing towards weak autocracy or otherwise cripple the state
\end{itemize}


\subsection{Important Specific Factual Claims}

\begin{itemize}[leftmargin=*, itemsep=4pt]
    \item Particular focus on the rise of the British Conservatives as the motivating success, as elites shifted to programmatic party politics
    \item Examines Germany as the failure case, both in Imperial and especially in Weimar Germany, where elites were excluded from electoral politics and instead had to pursue nondemocratic means of political engagement
    \item 
\end{itemize}


\subsection{Connections}

\begin{itemize}[leftmargin=*, itemsep=4pt]
    \item Provides an alternative explanation to Moore 1966: it is not how coalitions of classes formed or how wealth was distributed across the economy that determined the outcome of democracy, but rather the ability of the elite to organize into a mass party. (Also a broader dispute with Moore in that the motivating actor is the elites rather than the peasants)
    \item Huntington \& Dahl: super strong connection here, primarily regarding the ability of political institutions to appropriatey facilitate inclusive political participation. Ziblatt's claim is ultimately that without well-developed Conservative parties, elites were left without democratic influence and instead turned to anti-democratic means; conversely, when elite participation in democracy could be guaranteed by means of these parties, democracy was much more likely to succeed and be sustained
\end{itemize}


\subsection{My Critiques}

\begin{itemize}[leftmargin=*, itemsep=4pt]
    \item
\end{itemize}

\end{document}
